\documentclass[10pt,a4paper]{article}
\newcommand{\runningtitle}{Construct validation of single-cell topic loadings}
\input{layout}
\hypersetup{pdftitle={Mixture weight and programme strength: construct validation of single-cell topic loadings in diabetic foot skin}}
\begin{document}
\papertitle{Mixture weight and programme strength: construct validation of single-cell topic loadings in diabetic foot skin}
\section*{Abstract}
\noindent A higher mean single-cell programme score can reflect either stronger expression within cells or more cells occupying a high state. We examined this distinction before interpreting a fibroblast-associated topic as a diabetic-foot-ulcer (DFU) healing marker.
A frozen 15-topic model represented 76,461 cells from 25 selected foot-skin samples. We assessed distributional semantics, technical artifacts, analysis-constant sensitivity, negative controls and external projections. An author-derived mapping collapsed 14 DFU samples to 11 patients; anatomical controls and four representation benchmarks used the corresponding analysis units.
Across 25 samples, mean fibroblast loading correlated with the threshold-defined high-state fraction at $\rho=0.933$ (sample-bootstrap 95\% interval 0.789--0.983). The highest fraction, 0.841, occurred in healthy non-diabetic skin. Depth matching preserved 99.5\% of state labels, and exclusion of cells carrying immune transcripts preserved sample weights ($\rho=0.989$). Semantic associations persisted over tested thresholds and topic numbers; the healthy-skin counterexample failed at $K=10$, and healing-arm significance was threshold-sensitive. The corrected all-cell mean-loading contrast in 11 patients was 0.0280 (permutation $p=0.2719$; bootstrap interval $[-0.0017,0.0631]$). The separate 9-versus-5 high-state-fraction design had 24.9\% simulated power at its observed effect. Patient-level topic, module, PCA and NMF AUCs were 0.714, 0.786, 0.500 and 0.125; these exploratory estimates do not establish clinical superiority. Paired anatomy remained detectable in 10 patients. External cohorts did not supply independent patient-level healing validation.
The mean loading carries information about the prevalence of high-scoring cells, while state-specific means also vary. Its correlation with the threshold fraction does not prove invariant within-state intensity, and the tested fibroblast pattern remains an unvalidated prognostic marker. Patient-level validation requires independently defined outcomes and a clinically justified effect margin.
\par\smallskip\noindent Keywords: topic model; cell-state abundance; construct validity; DFU; CHI3L1; TIMP1; patient-level inference.

\section{Introduction}
Diabetic foot ulcers arise in tissue where inflammation, matrix remodeling and vascular impairment can persist together. Even after closure, recurrence remains common \citep{armstrong2017,wilkinson2020}. Single-cell studies have described fibroblast-associated programmes that differ between healing groups \citep{theocharidis2022,geo_gse165816}. Such observations make a programme score a plausible candidate measurement, but they leave a basic question unresolved: what changes when its average is higher in one biopsy than another?

There are at least two contributions. Cells already expressing a programme may express it more strongly, or the biopsy may contain a larger fraction of cells with high scores. Both can occur together. Topic models and gene-expression decompositions provide useful representations of heterogeneous cells \citep{blei2003,kotliar2019}, and neural latent-variable models can summarize expression for subsequent comparisons \citep{lopez2018}. Their coefficients nevertheless acquire biological meaning through the measurement and sampling scheme. A simplex constraint does not make a fitted coefficient an independently validated cell-state probability, and compositional changes can affect a tissue summary without an equivalent change within each cell \citep{aitchison1980,buettner2021}.

This distinction is difficult to resolve from a correlation between a mean and a threshold fraction: both are functions of the same cell scores. Distributional descriptions therefore need to be accompanied by an explicit decomposition of the mean and an examination of the terms treated as approximately constant. Technical explanations require separate attention. Dissociation-induced expression and ambient RNA can alter recovered profiles \citep{vandenbrink2017,young2020}, while normalization and cell selection affect which biological contrasts remain visible \citep{heumos2023}.

The unit of clinical inference introduces another constraint. Repeated biopsies and thousands of cells from a patient do not supply independent healing outcomes. Sample-aware methods address this dependence when testing expression or abundance changes \citep{crowell2020,squair2021,zimmerman2021,lazic2010}. Prognostic evaluation additionally requires a defined outcome, a relevant validation population and uncertainty assessed at the patient level \citep{collins2015}. Examining several readouts in a small discovery cohort can inform a later study, but selecting among them on that same evidence risks optimistic performance estimates \citep{cawley2010}. Similarly, a nonsignificant difference is not evidence of equivalence without a justified margin \citep{schuirmann1987,lakens2017}.

We examine a frozen fibroblast-associated topic in public diabetic-foot skin data, first as a tissue measurement and then as an exploratory healing readout. The analysis distinguishes mean loading within fibroblasts, the fraction above a fitted threshold and the mean over all recovered cells. It combines distributional and technical checks with author-mapped patient comparisons and paired anatomical controls. External wound datasets extend the tissue contexts in which the frozen score can be examined, while their available metadata determine whether they can address clinical validation.

\section{Results}
\subsection{Observed tissue composition provides the context for a fibroblast readout}

The study material comprises dissociated foot-skin specimens from diabetic-ulcer and non-diabetic groups, rather than longitudinally tracked cells (Figure~\ref{fig:p1f1}A\label{call:p1f1}). The 25 specimens represent 7 healed, 4 non-healed and 9 non-diabetic patients; cohort roles and analysis units are specified in Table~\ref{tab:p1cohorts}. The schematic in A separates model training from frozen inference, then branches into the fibroblast-restricted mean/fraction and the primary all-cell patient contrast. These are different estimands, not successive processing stages. The encoder and decoder in Equations~\eqref{eq:p1encoder} and~\eqref{eq:p1decoder} define the measured coordinate; the centered-log-ratio display in Equation~\eqref{eq:p1display} changes only its visualization. Mapping the saved expression coordinates makes the recovered cellular populations visible before comparing their statistical summaries. The display includes stromal, epithelial, vascular and immune compartments among 76,461 cells; fibroblasts account for 19,410 cells, or 25.4\% of the retained population (Figure~\ref{fig:p1f1}B). Coloring the same coordinates by the frozen fibroblast-associated loading shows how that continuous measurement relates to the operational lineage assignments (Figure~\ref{fig:p1f1}C). No new cell type or trajectory is defined by separation in this two-dimensional display.
\paperfigure{figure1_workflow}{fig:p1f1}{Observed study material and expression context. (A) Conceptual skin anatomy and sampling, not a tissue image or inferred spatial map. (B) A uniform 2,000-cell subset of the 76,461-cell UMAP, colored by marker-derived lineage. Fibro., fibroblast; EC, endothelial cell; mural, pericyte/smooth muscle. (C) The identical cells colored by frozen programme loading; the map and colors do not define new cell types. (D) Raw expression across all retained cells: dot size is the detection fraction and color the gene-wise scaled mean log-normalized expression. Annotation markers provide descriptive, not independent, support for labels. (E) All 19,410 fibroblast loadings in fixed-width histograms, with the existing threshold marked; healthy participants have no healing outcome. (F) Recovered-cell composition in each of 25 specimens, grouped by source outcome or non-diabetic status. Colors match B; columns are specimens, including repeated patients. Inference uses the relevant full data and patient units, not the display subsample.}


Raw-count reconstruction reproduced all saved lineage assignments and their sample order. The measured marker profiles provide an expression-level view of those assignments and of the programme-associated genes (Figure~\ref{fig:p1f1}D). Because some displayed markers also contributed to annotation, this is descriptive support rather than an independent validation of cell identity. Loading distributions within fibroblasts retain both low and high values across the clinical groups (Figure~\ref{fig:p1f1}E), while individual specimens differ in their recovered cell composition (Figure~\ref{fig:p1f1}F). Consequently, a pooled expression map, a cell fraction and a patient outcome contrast describe different aspects of the material. The mapped patient readouts are retained in Table~\ref{tab:p1biologicalpatients}.
\FloatBarrier

\subsection{The frozen representation is structured, but its principal loading is a mixture coordinate}
The discovery representation retained 76,461 quality-controlled cells from 25 selected foot-skin samples in the original 54-sample series. Beta perplexity was 243, and the manifold audit passed its recorded integrity checks. The connectivity statistic in Equation~\eqref{eq:p1reach} does not establish local geometric fidelity. The fitted loss in Equation~\eqref{eq:p1loss} uses Gaussian sampling and regularization as specified in Equation~\eqref{eq:p1gaussian}; deterministic readouts use the encoder mean, not a newly sampled latent state. Fibroblast topic was enriched for TIMP1, COL3A1, DCN, FN1, LUM and CHI3L1, with immediate-early and oxidative-stress genes also present among the top-ranked terms. The gene list therefore supports a fibroblast-associated inflammatory/stress description; it does not by itself establish a unique cell state or a healing mechanism.

Applying the mean--fraction identity in Equation~\eqref{eq:p1identity} to all 19,410 discovery fibroblasts clarified the remaining within-bin contribution. Empirical low-bin means ranged from 0.0120 to 0.1021, and nonempty high-bin means from 0.2264 to 0.5778. Nine of 25 specimens had an empty high bin. The approximation residual defined in Equation~\eqref{eq:p1residual} had RMS of 0.01575 and a residual range of $-0.02434$ to $0.04377$. The exact mean identity held to below $10^{-15}$ in every specimen (Table~\ref{tab:p1math}). The constant-state approximation is therefore explicit and imperfect; the correlation cannot establish invariant programme intensity.

A historical exploratory scan of 31 covariates across 25 specimens yielded 14 numerical BH q values below 0.05 under a naive independent-row analysis. Repeated patients invalidate that row-independence premise, so these are not 14 false-discovery-rate-controlled findings. The largest positive association was with the assigned B/plasma-cell fraction ($\rho=0.828$); CFD/APOD-rich topic mean loading was negatively associated ($\rho=-0.766$). Fibroblast PTPRC positivity and median UMI count were also associated ($\rho=0.727$ and 0.663). The complete scan, including non-significant covariates, is in Table~\ref{tab:p1covariates}. These marginal, sample-unit correlations do not establish causal drivers or independence from tissue composition; the positive depth association remains relevant despite the successful count-matching control.

\Needspace{65pt}
The hard-threshold fraction in Equation~\eqref{eq:p1fraction} uses a pooled two-Gaussian fit for a descriptive threshold; this fit did not test a continuous-shift model against a biological mixture model. High-loading samples retained a low-state tail overlapping the low-loading samples. Across the 25 analyzable samples, the sample mean and hard-threshold high-state fraction had Spearman $\rho=0.933$ (Figure~\ref{fig:p1f2}A\label{call:p1f2}). Quantile-based sample categories summarize the coexistence of low and high loadings without assigning a biological mechanism to those categories (Figure~\ref{fig:p1f2}B). The historical specimen-bootstrap interval of 0.789--0.983 and specimen-omission range of 0.924--0.957 are retained only as row-level diagnostics in Table~\ref{tab:p1semantics}, not patient-level uncertainty. The permutation null was centred near zero (Table~\ref{tab:p1semantics}). This is the central measurement result: in this dataset, the sample average is strongly associated with the high-state fraction. The two summaries are calculated from the same cells, so their correlation alone does not prove that component intensities are constant. The result does not say that high-state cells have no intrinsic biology. Together with the low-state tail, it shows that the mean contains information about high-score prevalence; it does not identify the relative biological contributions of prevalence and within-state expression. Patient aggregation gave a correlation of 0.937 across 20 patients, with a 10,000-draw patient-bootstrap 95\% interval of 0.776--0.990. Omitting one patient at a time gave correlations of 0.927--0.964 (Figure~\ref{fig:p1f2}C; Table~\ref{tab:p1robustness}). These descriptive patient-resampling summaries replace specimen resampling as the primary uncertainty display, without adding independent observations.

\paperfigure{figure2_mixture_semantics}{fig:p1f2}{Association between mean topic loading within fibroblasts and the high-state fraction. (A) Each point is one of 25 GSE165816 specimens: 9 healing (green circles), 5 non-healing (orange triangles) and 11 healthy non-diabetic (purple squares). The high-state fraction uses the pooled cut 0.1842958; Spearman ρ=0.933. (B) Sample-shape counts use within-sample loading quantiles: pure low requires the 90th percentile at most 0.15, pure high the 10th percentile at least 0.10, and mixed the 10th percentile at most 0.05 and 90th percentile at least 0.40; other samples form the remaining category. (C) Unlike the specimen-level panels A and B, uncertainty uses 20 author-mapped patients after pooling repeated specimens by fibroblast count. The dot is the patient correlation 0.937, the blue segment is the descriptive 10,000-draw patient-bootstrap 95\% interval (0.776--0.990), and the orange segment is the leave-one-patient-out range (0.927--0.964). Both correlated summaries use the same cells; the association does not by itself prove stable within-state intensity or a healing effect.}


\subsection{The operational high-state pattern survives selected technical controls}

The operational high-state pattern survived the originally specified artifact challenges, while the expanded controls identify measurement dependence. The dissociation stress standardized mean difference was +0.246, below the rejection threshold. PTPRC positivity was 0.98\% in high-state cells and 0.62\% in low-state cells, a rate ratio of 1.60 (Figure~\ref{fig:p1f3}A\label{call:p1f3}). Exact panel-count matching produced 99.5\% cell-label agreement and a specimen high-state-fraction correlation of 0.984 among 14,561 higher-depth retained fibroblasts under fixed lineage labels (Figure~\ref{fig:p1f3}B). This does not test the 4,849 excluded fibroblasts or rerun the quality-control and lineage-assignment pipeline after count reduction. The result therefore supports conditional score stability, not depth invariance of the entire tissue measurement. PTPRC positivity is only a proxy and does not exclude all doublets. A raw-count structureless control preserved every gene\textquotesingle{}s value multiset but changed the topic concentration substantially (perplexity 3.44 to 11.23 of 15), showing sensitivity to the joint perturbation of co-expression and row-wise library-size and detection-rate variation. The adaptive marker reference changed 9.21\% of cell labels and shifted the maximum specimen mean by 0.0371 and high-state fraction by 0.0695; there is no independent identity truth set to select between gates (Table~\ref{tab:p1gatecontrol}).

\paperfigure{figure3_artifact_controls}{fig:p1f3}{Technical and immune-transcript exclusion controls in the discovery cohort. (A) Dissociation standardized mean difference (SMD, 0.246) and the high-state/low-state PTPRC-positive rate ratio (1.60) are divided by rejection bounds 0.8 and 2.0; the dashed line is one. The rounded state cut is 0.184. (B) Each point is one of 25 samples after restricting to 14,561 fibroblasts eligible for panel-count matching to 986 counts per cell. Reprojection preserves sample fractions (ρ=0.984) and 99.5\% of cell labels; the dashed line is identity. (C) Bars are BIC for one Gaussian minus BIC for two Gaussians fitted to sample high-state fractions, using 24 matched samples before and after exclusion. Positive values favour two components. Exclusion retains 11,262 of 19,410 fibroblasts overall, and matched sample fractions correlate at ρ=0.989. This control measures invariance to the exclusion rule; it does not identify every source of ambient RNA or tissue-composition dependence.}


The full released-count audit reproduced the 81,602 pre-QC cells, 76,461 baseline survivors and all 19,410 fibroblast assignments, with a maximum deterministic-coordinate discrepancy below $4.77\times10^{-7}$. At 75\%, 50\%, 25\% and 10\% count retention, the three seeded runs retained 76,375--76,393, 75,724--75,738, 72,145--72,197 and 61,366--61,397 cells, respectively. Among cells surviving both QC passes, 3.18--3.30\%, 6.23--6.34\%, 10.29--10.48\% and 15.31--15.40\% changed lineage labels. At 10\%, only 80.11--80.15\% of original survivors remained, while 111--115 original failures newly passed QC. The gate assigned a winning label to 92 baseline QC survivors with no detected annotation marker; this increased to 1,410--1,462 at 10\% depth (Table~\ref{tab:p1fulldepth}). These labels are not identity truth. At 10\%, score rank correlation was about 0.910 on original retained cells with estimable scores, versus about 0.947 among QC survivors, exposing selection in a survivor-only stability claim. The full-pipeline all-cell patient contrast ranged from 0.0288 to 0.0296, with exact probabilities 109/330--114/330, versus baseline 90/330. All four readouts and rates are retained in Table~\ref{tab:p1depthoutcomes}. Persistence of a group contrast therefore coexists with substantial changes to the cells and labels being summarized.

Removing cells carrying any of the prespecified immune transcripts left the specimen high-state fractions nearly unchanged ($\rho=0.989$) and increased the BIC support for two components fitted to the sample fractions (Figure~\ref{fig:p1f3}C). This supports stability under the specified immune-transcript exclusion rule; it does not identify or exclude ambient RNA contamination. The placebo analysis was deliberately retained: unrelated tissue covariates flattened residual bimodality as well as immune covariates. At 24 matched samples, the data cannot determine whether the state is statistically separable from every tissue-composition axis. The matched exclusion comparison contains 24 samples, while each residual regression uses 25 samples. The immune and placebo regressions explained 49.5\% and 35.5\% of the variance, respectively, and both residual BIC differences favoured one component (Table~\ref{tab:p1ambient}). A clean technical control is not evidence that the state is a healing marker.

\subsection{High-state abundance does not establish healing specificity}

The lineage marker list shared no genes with the specified top-gene list used to describe the tested topic. This limited non-overlap does not make the gate independent of the full-panel learned representation. Fibroblast composition itself differed between the healing and non-healing sample groups only descriptively (ratio 1.42, $p=0.438$). Within the fibroblast gate, the discovery contrast was 2.89-fold, with a bootstrap 95\% interval of 1.06--13.97 and sample-level $p=0.04196$. This value is retained as a discovery observation, not as an independent patient validation, because the 14 samples represent 11 patients and include repeated same-wound, same-day specimens. The complete lineage and composition contrasts are in Table~\ref{tab:p1lineage}.

The high-state pattern was not restricted to samples labelled as healing. Four of nine healing samples and none of five non-healing samples were classified as high-mode samples, Fisher exact $p=0.2208$. The most extreme high-state fraction was a healthy non-diabetic foot-skin sample, whose high-state fraction was 0.841 (Figure~\ref{fig:p1f4}A\label{call:p1f4}). That observation is an existence counterexample to the statement that the state is uniquely a marker of a healing DFU. It neither identifies the tissue process responsible nor refutes a possible probabilistic association with healing.

\paperfigure{figure4_outcome_counterexample}{fig:p1f4}{State abundance, analysis unit and design power. (A) Symbols show high-state fractions for 9 healing, 5 non-healing and 11 healthy specimens; horizontal segments mark arm means. Green circles, orange triangles and purple squares identify these groups. The largest fraction is 0.841 in healthy non-diabetic foot skin, showing that the state is not exclusive to healing ulcers. (B) Dots and within-arm bootstrap 95\% intervals show the healing-minus-non-healing difference in all-cell mean fibroblast-associated topic loading for 14 specimens and 11 mapped patients. The patient analysis resamples 7 healing and 4 non-healing patients in 30,000 draws. Its readout differs from the fibroblast high-state fraction in A; the vertical dashed line is zero. (C) Gaussian location-shift simulation with a two-sided Mann--Whitney test estimates 24.9\% power for the separate 9-versus-5 high-state-fraction design, against an 80\% target. Its minimum detectable difference is 0.519. These are design sensitivities, not effect-exclusion bounds or clinical minimum important differences.}


For the separate all-cell mean-loading contrast in Equation~\eqref{eq:p1delta}, collapsing repeats to 11 patients gave a healing-minus-non-healing difference of 0.0293 (Figure~\ref{fig:p1f4}B). This is not a patient-level recalculation of the 2.89-fold fibroblast-restricted contrast. Exhaustive enumeration in Equation~\eqref{eq:p1exact} gave a two-sided probability of 0.272727 (90/330 allocations at least as extreme), versus 0.274924 under the retained add-one rule in Equation~\eqref{eq:p1permutation}; the bootstrap 95\% interval was {[}−0.0024, 0.0667{]}. The Welch standard error was 0.019446, and the data-derived equivalence infimum was 0.065375. The earlier 14-sample diagnostic contrast was 0.0457 with exhaustive $p=0.078921$ and an equivalence infimum of 0.079644; it is retained as a specimen-unit sensitivity analysis. The two-one-sided-test rule in Equation~\eqref{eq:p1tost} and boundary in Equation~\eqref{eq:p1margin} explain these data-derived sensitivity limits; neither margin is a clinical MCID. The corrected result is therefore a unit-aware discovery result with substantial uncertainty, not a definitive null test (Table~\ref{tab:p1exact}). Mapped scores are given in Table~\ref{tab:p1patients}; equal weighting of repeated specimen means gave 0.0297. Restricting the secondary specimen-label null to the 14 DFU specimens, rather than allowing healthy controls into pseudo-outcome groups, gave exact probabilities 0.0894, 0.1199 and 0.2697 for the fibroblast mean, high-state fraction and fibroblast composition, respectively; repeated specimens make these secondary quantities (Table~\ref{tab:p1dfunull}). A constructed continuous-score example produced a median mean--fraction Spearman correlation of 0.951, at least as large as the observed 0.933 in 76\% of draws, so the correlation cannot identify a discrete mixture without further assumptions. This construction is illustrative rather than a fitted null or calibrated probability. The full equivalence sensitivity grid is in Table~\ref{tab:p1equivalence}. Patient omission left the all-cell contrast positive, from 0.0143 to 0.0365 across 11 deletions (Table~\ref{tab:p1robustness}); this is an influence range, not a confidence interval or evidence of clinical discrimination.


\subsection{Retrospective design simulations do not resolve outcome uncertainty}

A retrospective simulation substituted the observed specimen-level high-state-fraction effect and pooled spread into a hypothetical 9-versus-5 independent-specimen design, giving 24.9\% estimated power. Under that simulation, the difference associated with 80\% power was 0.519 (Figure~\ref{fig:p1f4}C). The observed specimens include repeated patients, so this is neither the power of the corrected 7-versus-4 patient comparison nor independent evidence explaining its nonsignificance. These conditional design calculations cannot exclude a moderate or large healing association.

We also ran a design calculation for the next cohort, keeping it separate from the biological results. Under the deterministic specimen-derived arm spreads, hypothetical independent observations, 80\% target power and no missing outcomes, a candidate difference of 0.03 gave a total of 47, while a candidate equivalence margin of 0.05 gave 20. These spreads include repeated specimens and are not patient-population variance estimates. These values are design inputs rather than findings: the clinical MCID is still open, repeated samples increase the design effect, and the calculation cannot turn the discovery cohort into an adequately powered study. Tables~\ref{tab:p1design} and~\ref{tab:p1aucdesign} give the full difference, equivalence and exploratory AUC-gap design calculations.

\subsection{Fibroblast-restricted patient comparisons retain substantial uncertainty}

We next evaluated the two fibroblast-restricted readouts at the same independent patient unit used for the all-cell analysis. This closes the gap between the original sample-level lineage comparison and the mapped clinical outcomes. The denominator-specific pooling rules in Equation~\eqref{eq:p1lineagepatients} gave mean within-fibroblast loadings of 0.1729 in seven healed patients and 0.0837 in four non-healed patients. Their difference was 0.0892, with a 10,000-draw patient-bootstrap interval of $[-0.0486,0.2177]$ and an exhaustive two-sided permutation probability of 0.3121 (Figure~\ref{fig:p1f5}A\label{call:p1f5}). The corresponding high-state fractions were 0.3010 and 0.1496; the difference was 0.1514, with interval $[-0.1331,0.4153]$ and $p=0.3697$ (Figure~\ref{fig:p1f5}B). These comparisons use the existing gate and threshold, without searching for a cut that improves outcome separation.

\paperfigure{figure5_sensitivity_negative_controls}{fig:p1f5}{Threshold, capacity and negative-control analyses. (A,B) Spearman correlation of fibroblast mean loading with the high-state fraction across 25 samples and the two-sided Mann--Whitney healing-arm p value across 9 versus 5 samples, for all 18 cuts from 0.08 to 0.40. Vertical dashed lines mark the fitted cut; the horizontal line in B marks p=0.05. (C) Correlations for K=10, 15 and 20 after tracking the TIMP1-carrying topic; orange marks K=10, where the highest fraction is no longer in healthy skin. (D) The observed high-state-fraction arm difference is compared with the central 95\% of 20,000 label permutations (p=0.1358). (E) Gene-column shuffling of 6,000 cells raises mean cell perplexity from 3.68 to 14.43 out of 15, violating the concentration criterion. (F) The anatomical sample-level control compares 33 foot and 11 forearm specimens: observed p=0.0015, with 5.0\% of permuted comparisons below 0.05. This calibration is distinct from the paired 10-patient anatomy analysis and does not establish pipeline-wide error control.}


Fibroblast composition was evaluated separately, retaining all recovered cells as its denominator. The patient-level fractions were 0.2825 and 0.2118, a difference of 0.0708 with interval $[-0.0708,0.2146]$ and $p=0.4697$ (Figure~\ref{fig:p1f5}C). Benjamini--Hochberg adjustment across these three new readouts gave $q=0.4697$ for each. These are retrospective within-cohort comparisons; neither the positive point estimates nor their lack of significance establishes prognostic validity or equivalence.

Equal weighting of repeated specimen summaries within each patient gave differences of 0.0974, 0.1683 and 0.0679 for the programme mean, high-state fraction and fibroblast composition, respectively. All DFU specimens already contained at least 100 assigned fibroblasts, so imposing that requirement changed no outcome observation and supplies no independent robustness evidence. Omitting one patient at a time left the three differences positive, with ranges 0.0593--0.1445, 0.0850--0.2720 and 0.0279--0.1173, respectively. These influence ranges coexist with bootstrap intervals spanning zero; they are not confidence bounds. Table~\ref{tab:p1biologicalcontrasts} reports the complete weighting comparisons and exact probabilities.

The threshold-free patient distribution statistic in Equation~\eqref{eq:p1ecdf} was 0.3045, with exhaustive $p=0.1848$ over 330 patient-label allocations. Across the eight fixed cuts and the original midpoint, high-fraction differences remained positive (0.0570--0.2485); all unadjusted probabilities exceeded 0.22 and all support-wide max-statistic probabilities exceeded 0.31 (Figure~\ref{fig:p1f5}D; Table~\ref{tab:p1thresholdpatient}). Thus, the uncertain outcome association persists when the evaluation considers the whole loading distribution or multiple fixed thresholds, instead of relying only on the fitted midpoint.

Equalizing each specimen to 100 sampled fibroblasts left the programme-mean contrast positive in all 250 draws (0.0749--0.1233) and the high-state contrast positive in all draws (0.1098--0.2229); their smallest exact probabilities were 0.1273 and 0.1515 (Figure~\ref{fig:p1f5}E). Selecting one specimen per patient in all eight possible combinations gave programme-mean, high-state and composition ranges of 0.0777--0.1170, 0.1284--0.2081 and 0.0440--0.0918, respectively (Figure~\ref{fig:p1f5}F; Table~\ref{tab:p1samplingsensitivity}). These computations narrow the explanation that the observed direction is solely caused by unequal recovery or a particular repeated specimen, while retaining uncertainty about the healing association.


Direct comparison of the five repeated-specimen pairs, including three DFU patients, showed that within-specimen cell partitioning does not describe all between-specimen variation. In G7 and G8 from the same wound and day, fibroblast mean loading was 0.3504 versus 0.0921, an absolute difference of 0.2583. The separate within-specimen split-half 97.5th-percentile mean differences were 0.0553 and 0.0185; the full-distribution CDF supremum between specimens was 0.5789 (Table~\ref{tab:p1repeats}). This diagnostic leaves tissue microheterogeneity, preparation and annotation differences confounded. It does not estimate clinical repeatability or a patient-level significance probability.

\subsection{Method constants expose both stable and fragile claims}

The mean-versus-weight association remained positive across the cut grid, from $\rho=0.909$ to 0.948 (Figure~\ref{fig:p1f6}A\label{call:p1f6}), and across K=10, 15 and 20, from $\rho=0.924$ to 0.947 (Figure~\ref{fig:p1f6}C). The association therefore persisted over the tested cut and capacity choices.

\paperfigure{figure6_sensitivity_negative_controls}{fig:p1f6}{Threshold, capacity and negative-control analyses. (A) Each point gives the Spearman correlation between mean fibroblast loading and high-state fraction across 25 discovery specimens at one of 18 cuts from 0.08 to 0.40. The vertical dashed line marks the pooled fitted cut, 0.1842958. (B) Each point gives the two-sided Mann--Whitney healing-arm p value for 9 healing versus 5 non-healing specimens at the same cuts. The vertical line marks the fitted cut and the horizontal line marks p=0.05; 4 of 18 values fall below that level. (C) Correlations for K=10, 15 and 20 after tracking the TIMP1-carrying topic; orange marks K=10, where the highest fraction is no longer in healthy skin. (D) The observed high-state-fraction arm difference is compared with the central 95\% of 20,000 label permutations (p=0.1358). (E) Gene-column shuffling of 6,000 cells raises mean cell perplexity from 3.68 to 14.43 out of 15, violating the concentration criterion. (F) The anatomical sample-level control compares 33 foot and 11 forearm specimens: observed p=0.0015, with 5.0\% of permuted comparisons below 0.05. This calibration is distinct from the paired 10-patient anatomy analysis and does not establish pipeline-wide error control.}


The healthy-skin counterexample was more conditional. It held at all 18 tested cuts, but it failed at K=10 because the topic carrying TIMP1 belonged to a different IGFBP-rich programme. It held at K=15 and K=20 when the CHI3L1/COL3A1/DCN/FN1 state resolved as its own topic. The defensible wording is consequently ``the counterexample holds for the tested K=15 and K=20 fits but not the K=10 fit,'' not ``the counterexample is invariant to model capacity.''

The healing-arm p value moved between 0.042 and 0.062 over the cut grid, with only 4 of 18 cuts below 0.05 (Figure~\ref{fig:p1f6}B; Tables~\ref{tab:p1cuts} and~\ref{tab:p1capacity}). That sensitivity is a reason to avoid presenting the discovery p value as a stable healing signal. It also independently supports the decision to keep the healing claim unvalidated.

\subsection{Negative controls and representation comparisons constrain interpretation}

The DFU-only label-permutation calibration used the 9 healing and 5 non-healing specimens, excluding all 11 healthy controls before allocation. Its observed high-state-fraction difference was +0.2380, with a central 95\% null interval {[}−0.3167, +0.3222{]} and exact $p=0.1199$ over 2,002 allocations (Figure~\ref{fig:p1f6}D). Because this null relabels specimen outcomes and includes repeated patients, it is a secondary specimen-unit calibration rather than a patient-level clinical test. The observed contrast was not exceptional under this label-exchangeability null, and this single calibration does not prove pipeline-wide false-positive control. The raw-count structure control used a fixed 6,000-cell source-order subset and shuffled gene values across cells while preserving each gene\textquotesingle{}s observed value multiset; real data had perplexity 3.44/15, whereas the shuffled input had 11.23/15 and lower HHI (0.574 to 0.147). The shuffle therefore reduced coordinate concentration while also narrowing row-wise library-size and detection-rate variation; because it does not preserve row-wise library sizes, it is a co-expression-structure control rather than a complete technical null (Figure~\ref{fig:p1f6}E; Table~\ref{tab:p1rawcontrol}). In the anatomical permutation control, the observed p value was 0.0015 and 5.0\% of permuted comparisons fell below 0.05 (Figure~\ref{fig:p1f6}F). That fraction describes this permutation reference; it is not an independent repeated-dataset test of pipeline-wide type-I error.

The historical specimen-level scoring procedure used the pairwise AUC statistic in Equation~\eqref{eq:p1auc} and gave point AUCs of 0.822 for fibroblast-associated topic, 0.844 for the module score, 0.533 for PCA and 0.878 for NMF (Figure~\ref{fig:p1f7}A\label{call:p1f7}). No pairwise fixed-score bootstrap interval for an AUC difference excluded zero. The module score alone separated its procedure-specific permutation null at $p=0.0424$; fibroblast-associated topic and NMF gave $p=0.0604$ and $p=0.0597$, while PCA gave $p=0.4527$. Null means also differed from 0.5 (fibroblast-associated topic, 0.562; NMF, 0.553). These historical outputs pool raw coordinates from different fitted folds and include repeated patients; a permutation calculation does not make those coordinates comparable. They document the earlier procedure rather than supply evidence for a representation ranking (Tables~\ref{tab:p1sampleauc}, \ref{tab:p1correlations} and~\ref{tab:p1aucdifferences}).

The sample-level readouts had pairwise Spearman correlations from 0.174 to 0.710 (Figure~\ref{fig:p1f7}B). Their relatedness does not determine the ordering of their discrimination estimates.

\paperfigure{figure7_representation_inference}{fig:p1f7}{Historical specimen-unit scoring diagnostics in 14 DFU specimens (9 healing, 5 non-healing). (A) Points are leave-one-sample-out areas under the ROC curve (AUCs); horizontal bars are percentile 95\% intervals from 4,000 arm-stratified bootstrap draws of fixed out-of-fold scores, with common resampling indices across methods. Orange ticks mark the procedure-specific permutation-null means after rerunning outcome-dependent fitting and orientation. The dashed line at AUC=0.5 is a visual reference, not the fitted permutation null. Fibroblast topic denotes the frozen all-cell mean, Module the specified fibro-inflammatory score, PCA principal component analysis, and NMF nonnegative matrix factorization. Raw PCA/NMF coordinates can differ across fitted folds, so neither their pooled AUCs nor the displayed bootstrap contrasts support a representation ranking. (B) Cells contain Spearman correlations of these 14 specimen-level readouts; the colour scale and printed values encode the same statistic. Repeated patients and cross-fold comparability limit these historical diagnostics.}


\subsection{Same-model patient-pair discrimination and paired anatomical sensitivity}

The revised held-pair statistic in Equation~\eqref{eq:p1pairauc} evaluated all 28 healing/non-healing pairs from the 11 patients, fitting each scoring rule on the other nine patients. Within-model pair discrimination was 0.714 for the frozen topic (20 wins, no ties), 0.786 for the module (22 wins, no ties), 0.500 for PCA (14 wins, no ties) and 0.232 for NMF (five wins and three half-credit ties). These descriptive estimates use comparable scores within each pair and do not pool raw coordinates across fitted models (Table~\ref{tab:p1pairauc}; Figure~\ref{fig:p1f8}A\label{call:p1f8}). Changing held-patient expression or labels did not change fold preprocessing, component selection or orientation in the isolation tests; positive component rescaling also preserved NMF selection and pair order. The comparisons nevertheless reuse 11 discovery patients and a discovery-trained panel and encoder, and establish neither external validation nor a population representation ranking.

The earlier leave-one-patient-out procedure remains separately documented: topic AUC was 0.714 with add-one permutation $p=0.2991$, module AUC was 0.786 with $p=0.166$, PCA was 0.500 and NMF was 0.125. Its fixed-score bootstrap intervals and pairwise contrasts are historical procedure diagnostics, not uncertainty estimates for the revised held-pair estimand (Tables~\ref{tab:p1patientauc} and~\ref{tab:p1aucdifferences}). The change in the NMF estimate reflects the revised fold, projection and component-selection procedure; it is not an isolated estimate of a fold-scale effect.

\paperfigure{figure8_patient_unit_anatomy}{fig:p1f8}{Patient-unit discrimination and paired anatomical sensitivity. (A) Points are leave-one-patient-out AUCs for 7 healing and 4 non-healing patients after cell-count-weighted collapse of repeated specimens. Bars are 95\% intervals from 4,000 arm-stratified patient bootstrap draws of fixed out-of-fold scores; the dashed line marks AUC=0.5. Fibroblast topic, Module, PCA and NMF denote the frozen topic mean, fibro-inflammatory module, principal component and nonnegative matrix factorization readouts. (B) The point is the mean all-cell fibroblast-associated topic loading difference, foot minus forearm, across 10 paired patients; the bar is a 4,000-draw paired-patient bootstrap 95\% interval and the horizontal line is zero. The exhaustive sign-flip test with add-one correction gives p=0.0107 and Benjamini--Hochberg q=0.0402 across 15 topics. Both panels use the discovery source and are not independent clinical validation.}


The paired anatomical control confirmed that the frozen encoder responds to a known tissue difference. Across 10 paired patients, the fibroblast-associated topic increased in foot relative to forearm skin by +0.0476 (exhaustive sign-flip $p=0.0107$, BH $q=0.0402$); the fibroblast-restricted contrast was +0.1398 ($p=0.0146$). Four of 15 topics remained significant after BH correction (Figure~\ref{fig:p1f8}B; Table~\ref{tab:p1anatomy}). The earlier 3-of-15 sample-level positive-control result is directionally consistent, but the paired analysis is the more appropriate patient-unit sensitivity check. These controls make a complete assay failure less likely, but they do not turn the underpowered DFU comparison into a proof of no association.

\subsection{External projections remain construct audits}

The external immune-axis analysis used mutually exclusive lineage compartments. GSE223964 and GSE245703 gave immune-axis correlations of −0.024 and +0.846, respectively (Figure~\ref{fig:p1f9}A\label{call:p1f9}). Within the immune compartment, B/plasma-cell share was weakly related to the topic signal (−0.071 and +0.077), while correlations with median detected genes per cell were negative in both datasets (−0.429 and −0.427; Figure~\ref{fig:p1f9}B). The signs are inconsistent for a generalized immune-axis effect, and neither dataset has an authoritative patient-level healing endpoint (Table~\ref{tab:p1external}).

\input{figure_captions/f9}

A separate projection into GSE231643 provides a weaker, label-only concordance of direction: the within-fibroblast ratio falls from 2.89 in the discovery analysis to 1.66 (sample-bootstrap 95\% interval 1.01--6.10; two-sided Mann--Whitney $p=0.1429$; Table~\ref{tab:p1lineage}). The separate GSE241132 acute-wound cohort gave a broad immune-axis rho of −0.070 and an author-annotated B/plasma-cell rho of +0.210 across three donors; both values are descriptive and were not combined with the DFU outcome analyses. GSE268834 supplied another descriptive projection: the topic-associated fibroblast signal correlated with exclusive immune fraction at +0.714, and the diabetic-minus-non-diabetic fibroblast-topic contrast was +0.0464 with $p=0.2812$ (Figure~\ref{fig:p1f9}C). This is compatible with an immune/tissue-context hypothesis but cannot distinguish it from other between-sample differences.

GSE248247 was a deliberately low-yield stress test of the frozen projection. It contained 561 raw cells and 358 quality-controlled cells, with 6,994 of 7,002 panel genes covered in each sample. The fibroblast topic contrast between one diabetic and one non-diabetic sample was −0.0204 with descriptive permutation $p=1.0$ (Figure~\ref{fig:p1f9}C). All correlations were non-estimable because n=2. The exclusive compartment sum was exactly one in both samples, demonstrating that the lineage gate was non-overlapping; the dataset cannot validate healing or support pooled patient inference.

\FloatBarrier
\section{Discussion}
Mean fibroblast loading and the fraction of high-scoring fibroblasts were closely associated across the discovery specimens, and the association persisted after patient aggregation. This establishes that the mean conveys information about the prevalence of high scores among recovered fibroblasts. It does not make the measurements interchangeable. Low- and high-bin means varied across specimens, and some specimens contained no high-bin cells. An approximation that replaced these means with pooled values left a measurable residual. The defensible interpretation therefore includes both prevalence and within-bin score variation, conditional on the representation, gate and threshold used to define the measurement \citep{altman2006,buettner2021}.

The patient-unit extension connects the fibroblast interpretation directly to the available healing categories. Averaging programme scores within fibroblasts, counting high-scoring fibroblasts and counting fibroblasts among all cells all gave higher point estimates in healed patients, but each comparison retained considerable uncertainty. Reweighting repeated biopsies did not remove that uncertainty. A patient-level comparison of the entire loading distribution also retained uncertainty, while equal-cell sampling and every one-specimen-per-patient choice preserved the direction of the point contrast. These checks address selected measurement and sampling choices within the discovery material; they do not validate the readout in a new cohort. The appropriate implication is that the observed group direction warrants prospective evaluation of a fixed measurement, not that a threshold has been clinically validated or that a null biological effect has been demonstrated. Correcting the unit of analysis is necessary even when a sample-level contrast appears more favorable \citep{squair2021,zimmerman2021}.

The different uncertainty summaries answer different questions. Patient bootstrap intervals resample the observed outcome arms and inherit the limitations of their small empirical distributions \citep{efron1993}. Exhaustive label enumeration includes the observed allocation, whereas an add-one correction is commonly motivated for randomly sampled permutations \citep{phipson2010,hemerik2018}. Retaining the earlier convention alongside the exhaustive fraction makes that arithmetic transparent without making observational labels randomized. Likewise, Benjamini--Hochberg adjustment concerns a stated testing family under its dependence assumptions \citep{benjamini1995}; it cannot repair repeated specimens treated as independent patients. The retrospective power simulation substitutes the discovery effect and specimen spread into a hypothetical sampling model. It supplies no independent evidence about whether the uncertain patient association is real, absent or clinically important.

The cell atlas and expression profiles also clarify what is still missing experimentally. A loading enriched in stromal cells is not a direct count of an anatomically defined fibroblast population. Independent spatial or multiplex protein measurements in mapped biopsies would be needed to localize high-scoring cells and test whether their abundance changes within comparable tissue compartments \citep{foster2021,philippeos2018}. Perturbing a selected programme in an appropriate wound system could then ask whether it contributes to repair, rather than merely accompanying it \citep{mascharak2021}. These experiments address localization and function; neither can be supplied by adding more resamples to the present cell matrix.

This distinction has biological consequences. A high mean could accompany recruitment or recovery of a particular fibroblast population, increased expression within an existing population, or changes in several populations at once. Studies of dermal lineages and cross-tissue fibroblast organization show why these possibilities cannot be resolved by a small marker list alone \citep{driskell2013,muhl2020,buechler2021}. The enrichment of matrix and inflammatory genes supplies a useful description of the loading, but does not identify a unique lineage, a stable attractor or a causal repair mechanism. The broader separation between cell identity and cellular activity remains relevant when a continuous transcriptional programme is summarized at the tissue level \citep{kotliar2019}.

The threshold should be understood as an operational convention. It permits a reproducible high-state fraction and makes the denominator explicit, while assigning nearby scores on opposite sides of the cut to different bins. Information within those bins remains in the continuous score. This is one reason to retain both summaries rather than interpreting the binary fraction as a complete description of the programme \citep{altman2006}. The sensitivity to topic capacity also shows that a gene-associated coordinate need not have identical meaning across independently fitted representations. Neighborhood abundance methods could provide a complementary analysis of continuous state space, but they would introduce their own representation, neighborhood and replication choices \citep{dann2022}. They cannot retrospectively turn a descriptive threshold into proof of a discrete biological state.

The technical controls provide evidence about selected perturbations of this operational measurement. Most state labels survived count matching under fixed lineage assignments, and specimen fractions were stable after specified immune-transcript exclusions. Rerunning QC and the frozen gate on all released count columns reveals a different result: low depth removes many original cells, changes lineage labels among survivors and increases forced assignments without detected marker support. Stability conditional on a fixed or surviving population is therefore not stability of the whole tissue measurement. A relatively stable patient contrast can coexist with changed population membership and should not obscure that sensitivity. Neither stress test establishes cell-identity truth or excludes ambient RNA and doublets. Dedicated contamination and doublet procedures address different aspects of those processes \citep{young2020,wolock2019}. Dissociation responses can remain entangled with genuine injury responses \citep{vandenbrink2017}, and the similar attenuation of residual bimodality by immune and unrelated covariates limits attribution to one immune mechanism.

The difference between a tissue measurement and a healing predictor becomes clearer at the patient level. The healthy specimen with the highest fraction challenges disease exclusivity, while remaining compatible with a probabilistic association in a defined clinical population. Collapsing repeated specimens reduced the all-cell healing contrast and made its uncertainty explicit. This comparison concerns a different readout from the fibroblast-restricted fraction, so their effect estimates should not be treated as interchangeable measures of the same clinical association. More cells improve characterization of the sampled tissue but do not increase the number of independent outcomes \citep{squair2021,zimmerman2021}. The available healing categories also do not support a model of individual time to closure \citep{armstrong2017}.

The representation comparison remains an internal audit. Historical pooled-fold PCA and NMF coordinates are not directly comparable across fitted models, and their bootstrap contrasts cannot establish representation quality. The revised held-pair procedure uses the same fitted scoring rule for each healing/non-healing comparison and makes NMF component selection invariant to positive rescaling. It repairs that computational comparison without creating independent data: all pairs reuse the same 11 patients, and the frozen topic encoder and gene panel already used discovery expression. Overlapping training folds also share information \citep{bengio2004}. Exploring several thresholds, capacities and representations in the same cohort makes independent evaluation especially important, even when each individual comparison has a clearly specified computational rule \citep{cawley2010}. The uncertainty concerns the selection and transport of the readout as well as the arithmetic of a particular discrimination estimate.

Even reliable discrimination would leave other clinical questions unanswered. AUC describes ordering of outcomes, whereas calibration concerns whether predicted risks agree with observed risks in the population where the model is used \citep{steyerberg2010,vancalster2019}. The current loadings are expression summaries rather than calibrated clinical probabilities. Their value for a decision would also depend on the consequences of false reassurance and unnecessary intervention at clinically meaningful thresholds \citep{vickers2006}. This study does not fit such a decision model. Its contribution is to make the measurement and analysis unit sufficiently explicit that a later prognostic study can evaluate incremental value beyond routinely available clinical information.

The paired anatomical contrast demonstrates sensitivity to a known tissue difference, but it also emphasizes the importance of sampling context. Fibroblast organization varies across dermal compartments and with age \citep{driskell2013,soleboldo2020}. A detectable contrast between foot and forearm therefore supports responsiveness of the frozen representation without establishing specificity for wound healing. Spatial localization and an independently measured protein or transcript readout would help distinguish altered cellular abundance from altered activity within a compartment \citep{foster2021,kotliar2019}. Such measurements would address a different uncertainty from increasing the number of cells sequenced from the same biopsy.

External projection adds tissue contexts in which the frozen coordinate can be computed. It does not make the cohorts interchangeable. Different sampling designs, gene coverage, cell yields and compartment definitions can alter a summary before any difference in healing biology is considered. Benchmarking of single-cell integration demonstrates the importance of assessing technical alignment and conservation of biological variation together \citep{luecken2022}. Freezing the representation reduces one source of analytical flexibility, but does not ensure comparable measurement across source studies. The inconsistent immune associations and absence of independently adjudicated patient healing outcomes limit these cohorts to construct evaluation. The smallest projection, with two specimens, is chiefly a check of computational feasibility.

Future evaluation needs a defined outcome horizon, fixed sampling and readout, and a prespecified prognostic estimand. Prediction-model sample size also depends on outcome frequency and model complexity \citep{riley2019}. Equivalence requires a justified margin, not the smallest margin compatible with the observed data \citep{lakens2017}; a margin is not a prerequisite for every association or discrimination analysis. The present contribution is an executable measurement audit, not a new topic-model algorithm or biological law. It quantifies the limits of equating mean loading with prevalence, conditional stability with full-pipeline robustness, or specimen contrasts with patient outcomes. Computational reproducibility makes those limits inspectable; independent patients and orthogonal measurements are needed to establish clinical value.

\section{Materials and methods}
\subsection{Cohort selection and independent units}
GSE165816 comprises 54 original samples and 27 author-mapped participants. The frozen model used 25 selected foot-skin samples in the non-diabetic, DFU-healing and DFU-non-healing groups: 81,602 cells before cell quality control and 76,461 afterwards. Cells required at least 200 detected genes and a mitochondrial count fraction at most 0.20; genes required detection in at least 10 cells before panel selection. The final panel contained 7,002 genes, obtained from a nominal 7,000-gene dispersion selection with retained anchor genes. The discovery outcome subset contained 9 healing and 5 non-healing samples.

The clinical supplement maps those 14 outcome samples to 11 patients: 7 healing and 4 non-healing. G4/G23, G7/G8 and G33/G34 are repeated samples from the same wound on the same day. We combine each patient's retained cells before computing the primary patient score. The outcome is the study's healing category within a 12-week follow-up window, not a measured individual healing time. Series metadata describe a shared Emory 10x Chromium 3$'$ V2/V3 workflow and NovaSeq instrumentation; the absence of per-sample chemistry and flow-cell identifiers prevents run-specific batch adjustment.

\subsection{What the frozen neural topic model computes}
For nonnegative panel counts $x_{ig}$ with positive total, let $u_{ig}=x_{ig}/\sum_{h=1}^{G}x_{ih}$. The implemented model uses a diagonal-Gaussian encoder in unconstrained coordinates and a softmax decoder:
\begin{align}
q_\phi(z\mid u_i)&=\mathcal N\!\left(\mu_\phi(u_i),\operatorname{diag}\sigma_\phi^2(u_i)\right),\label{eq:p1encoder}\\
\theta_{ik}&=\frac{e^{z_{ik}}}{\sum_{\ell=0}^{K-1}e^{z_{i\ell}}},\qquad
\beta_{kg}=\frac{e^{b_{kg}}}{\sum_{h=1}^{G}e^{b_{kh}}},\qquad
\widehat u_{ig}=\sum_{k=0}^{K-1}\theta_{ik}\beta_{kg}.\label{eq:p1decoder}
\end{align}
Equations~\eqref{eq:p1encoder}--\eqref{eq:p1decoder} define the encoder, simplex coordinates and normalized topic-gene decoder used below. Thus $\theta_i\in\Delta^{K-1}=\{a\in\R^K:a_k\geq0,\sum_{k=0}^{K-1}a_k=1\}$ and every topic-gene row sums to one. The topic indices start at zero. A Gaussian mapped through softmax gives a logistic-normal representation. This differs from a Dirichlet prior, while retaining the mixture representation that motivates topic models. Training minimizes
\begin{equation}
\mathcal L(\phi,b)=\frac1N\sum_i\E_{q_\phi}\!\left[-\sum_g u_{ig}\log(\widehat u_{ig}+10^{-10})\right]
+0.01\,\frac1N\sum_i\operatorname{KL}\!\left(q_\phi(z\mid u_i)\,\|\,\mathcal N(0,I)\right).\label{eq:p1loss}
\end{equation}
Equation~\eqref{eq:p1loss} is the fitted objective: each cell contributes cross-entropy on its panel proportions, followed by a Gaussian regularizer. It is neither a depth-weighted count likelihood nor an unmodified evidence lower bound for such a likelihood. With encoder log-variance $\ell_{ik}=\log\sigma_{ik}^2$, the reparameterization and regularizer are
\begin{equation}
z_{ik}=\mu_{ik}+e^{\ell_{ik}/2}\epsilon_{ik},\quad \epsilon_i\sim\mathcal N(0,I),\qquad
\operatorname{KL}_i=\frac12\sum_k\left(\mu_{ik}^2+e^{\ell_{ik}}-1-\ell_{ik}\right).\label{eq:p1gaussian}
\end{equation}
Equation~\eqref{eq:p1gaussian} is the diagonal-Gaussian calculation used by the implementation; one noise draw per cell estimates the reconstruction expectation in each training batch. The encoder has widths 256 and 128, batch normalization, ReLU activations and dropout 0.1. The recorded fit used $K=15$, 40,000 Adam steps, learning rate $10^{-3}$ and batches of 512. Deterministic projection disables encoder dropout and uses $z=\mu_\phi(u)$ and $\theta=\operatorname{softmax}(\mu_\phi(u))$, which is not generally $\E_q[\operatorname{softmax}(z)]$. Missing panel genes in external projections are set to zero. Normalization and quality control can alter biological interpretations, so this definition is kept distinct from other variance-stabilizing methods.

Topic-gene perplexity is summarized as $\exp\{-K^{-1}\sum_{k,g}\beta_{kg}\log\beta_{kg}\}$; it measures concentration of decoder weights, not superiority over other representations. The latent coordinate $\mu\in\R^K$ is not the simplex coordinate $\theta$. Softmax is unchanged by adding a common scalar to all coordinates of a cell; reconstruction therefore does not identify that offset. The Gaussian penalty constrains the fitted coordinates but does not give them a unique biological interpretation. Topic weights are decoder coefficients, not independently calibrated probabilities of membership in a biological cell type.

\subsection{Threshold fractions and the exact mean decomposition}
Let $Y_{si}=\theta_{si,\star}$ for fibroblast-gated cell $i$ in sample $s$. The discovery lineage gate averages marker expression standardized across the selected discovery cells, then assigns the highest-scoring lineage. Its markers are disjoint from the tested topic-gene list by an explicit overlap check. The gate remains an expression-based operational definition; its reference population can change its assignments, and the frozen encoder alone does not freeze the complete tissue measurement. A pooled two-Gaussian fit provides a descriptive cut $c=0.1842958$, the midpoint of the fitted component means. Gaussian components approximate a bounded score distribution and are not assumed to be literal biological populations.
\begin{equation}
\bar Y_s=\frac1{n_s}\sum_iY_{si},\qquad
\widehat w_s(c)=\frac1{n_s}\sum_i\ind\{Y_{si}>c\}.\label{eq:p1fraction}
\end{equation}
Equation~\eqref{eq:p1fraction} defines both the sample mean and the reported hard-threshold fraction; the latter is not an average posterior responsibility of the Gaussian mixture. Define $\bar Y_{s,L}$ and $\bar Y_{s,H}$ as means below/at and above the cut. Then, with empty-component contributions interpreted as zero,
\begin{equation}
\bar Y_s=(1-\widehat w_s)\bar Y_{s,L}+\widehat w_s\bar Y_{s,H}.\label{eq:p1identity}
\end{equation}
Equation~\eqref{eq:p1identity} is exact. Only under an additional approximate component-stability assumption, $\bar Y_{s,L}\approx m_L$ and $\bar Y_{s,H}\approx m_H$, does it reduce to $\bar Y_s\approx m_L+(m_H-m_L)\widehat w_s$. Consequently, a large correlation between the two summaries cannot independently prove stable intensity: both summaries use the same cells. Persistence of overlapping low-state tails supplies complementary distributional evidence. The 10th percentile of a high-loading sample (0.029) overlaps the scale of the 90th percentile in a low-loading sample (0.119); this observation alone is not a formal test of multimodality.

To quantify the term omitted by a constant-state approximation, let $m_L,m_H$ be the cell-pooled means in the two threshold bins across the 19,410 discovery fibroblasts. They are empirical bin means, not the fitted Gaussian means. For each specimen,
\begin{align}
A_s&=(1-\widehat w_s)m_L+\widehat w_sm_H,\notag\\
R_s&=(1-\widehat w_s)(\bar Y_{s,L}-m_L)+\widehat w_s(\bar Y_{s,H}-m_H),\qquad \bar Y_s=A_s+R_s.\label{eq:p1residual}
\end{align}
Equation~\eqref{eq:p1residual} separates a specified constant-state approximation from its exact residual. An empty bin contributes zero and its undefined mean is omitted from the corresponding mean range. We summarize all 25 specimens using specimen-weighted residual RMS and bin-mean ranges. These are finite-data descriptions without a test of state invariance or a causal allocation of variance (Table~\ref{tab:p1math}).

The semantic statistic is $\rho_{\rm sem}=\operatorname{cor}_{\rm Spearman}(\bar Y_s,\widehat w_s)$ across 25 samples. A 10,000-draw sample-row percentile bootstrap, 10,000 row permutations and leave-one-sample-out recalculation characterize its stability. These do not create independent patients. The patient-collapsed correlation of 0.937 across 20 mapped patients and its descriptive patient-resampling interval are the primary uncertainty summary; the earlier specimen resampling is historical.

A separate patient-resampling analysis collapses the 25 specimen means and high-state fractions to 20 patients using fibroblast counts as within-patient weights. It draws 10,000 patient samples with replacement (seed 217) and reports the percentile interval for their Spearman correlation. Leave-one-patient-out analyses recompute this correlation across the remaining 19 patients and separately recompute the all-cell healing contrast across the remaining 10 outcome patients. The reported omission range is a sensitivity diagnostic, not an interval estimate (Table~\ref{tab:p1robustness}).

\subsection{Technical controls and constant sensitivity}
Dissociation-induced programmes and ambient transcripts motivate the technical controls. The stress comparison uses a standardized mean difference (SMD) with rejection bound 0.8; the doublet-oriented control compares PTPRC-positive rates with ratio bound 2.0. Exact panel-count matching reprojects downsampled counts, comparing cell labels and sample fractions. The historical 31-covariate screen retains its specimen-row asymptotic p values and BH transforms for auditability, not as valid patient-level significance or false-discovery-rate control: repeated patients violate the naive independent-row premise. Removing any cell carrying a specified immune transcript tests invariance to that exclusion rule. This cannot by itself identify ambient contamination or exclude every doublet. Residual mixture checks with immune covariates are compared with four unrelated placebo covariates; similarly flattened residuals limit claims of statistical separability.

Sensitivity fits use 18 cuts between 0.08 and 0.40 (including the rounded fitted cut) and $K\in\{10,15,20\}$. The TIMP1-carrying topic is tracked descriptively across capacities; topic labels are not assumed exchangeable or biologically identical across $K$. The anatomical positive-control calibration uses 20,000 shuffles. The raw-count structure control selects a fixed 6,000-cell subset by the recorded source order, then permutes each gene's observed values across cells with a seeded permutation; this preserves the marginal multiset for each gene but does not preserve row-wise library sizes. The resulting topic concentration, row-wise library-size and detection-rate summaries are compared with real input; no new topology audit is performed. The separate outcome-label calibration exhaustively enumerates all 2,002 allocations of the 9 healing and 5 non-healing specimens, retaining repeated specimens as a specimen-unit null. It is not a patient-level clinical test or proof of pipeline-wide error control. The DFU-only label null excludes all healthy controls before allocation and reports the 9/5 specimen result as secondary to the 7/4 patient test. For each specimen, a continuous-score construction applies the inverse-logit transform to its observed mean on the logit scale plus independent standard-normal noise, retaining its cell count; 200 seeded draws are made without a latent class or mixture fit. The constructed mean need not equal the observed mean. This is an illustrative non-identification example, not a fitted null model or a calibrated probability. These analyses are retrospective computational controls, not a registered clinical trial.

For marker-reference sensitivity, expression is normalized per cell to 10,000 counts and transformed by log1p. Each lineage score averages its marker-wise standardized values, followed by exclusive maximum-score assignment. The pooled marker means and SDs are frozen before applying the same reference to each specimen; this implementation check is compared with re-estimating means and SDs separately within each specimen. The topic encoder and state threshold remain fixed. Patient contrasts are recomputed after count-weighted pooling under each annotation reference, with the three readouts forming a separate BH family within each gate. This evaluates dependence on the annotation reference without an independent cell-identity truth set (Table~\ref{tab:p1gatecontrol}).

The exploratory covariate scan retains the original rounded rule $Y_{si}\geq0.184$, used by the technical-control analyses, rather than substituting the fitted cut 0.1842958. It correlates the sample high-state fraction with 14 other fibroblast topic means, 11 whole-sample lineage fractions, total and fibroblast cell counts, fibroblast dissociation score, median UMI count, PTPRC-positive fraction and the numeric G specimen label. Two-sided Spearman probabilities are adjusted across all 31 comparisons by Benjamini--Hochberg. The 25 rows are specimens, include repeated patients, and receive no clinical-outcome adjustment. The G label is an identifier covariate, not a verified sequencing batch or collection time. Topic coordinates and lineage fractions have compositional dependence.

\subsection{Full released-count depth sensitivity}
The full-pipeline stress test starts from all 81,602 released count columns in the 25 selected specimens, including cells failing the original analysis QC. These are processed count matrices, not pre-cell-calling droplets or sequencing reads. Baseline reconstruction must reproduce retained cell identities, lineage assignments and deterministic coordinates before perturbation. Each nonzero count is independently thinned with a binomial probability of 0.75, 0.50, 0.25 or 0.10, with three fixed computational seeds per rate. Every rate starts from the original counts; no cells or patients are resampled. QC is rerun using at least 200 detected genes and a mitochondrial fraction at most 0.20. The discovery marker reference, exclusive maximum-score gate, panel, encoder and high-state threshold remain unchanged; there is no new margin or outcome-guided label rule.

The full-pipeline population includes all current QC survivors and uses newly computed labels. Separate conditional branches retain original QC cells and labels, retain the QC intersection with original labels, or re-gate the same intersection. Lost and newly admitted cells, lineage transitions and fibroblast entry/exit are reported separately. Zero detected annotation markers and tied winning scores are diagnostics, not evidence of valid cell identity. A zero panel total makes a score missing, rather than assigning the encoder output of a zero vector. Readouts preserve their selected-cell denominators and are non-estimable if required scores are missing. Patient aggregation pools the appropriate all-cell or fibroblast counts across repeated specimens. The same seven healed and four non-healed patients supply 330 exhaustive label allocations; the three seed ranges describe perturbation sensitivity, not biological confidence intervals or new outcome tests. Uniform thinning does not model upstream cell calling, cell-type-specific capture, ambient RNA, doublets or independent tissue sampling.

\subsection{Patient contrast, permutation and equivalence sensitivity}
Let $\mathcal C_p$ contain all retained cells in the mapped DFU foot specimens of patient $p$. The primary corrected score is $S_p=|\mathcal C_p|^{-1}\sum_{i\in\mathcal C_p}\theta_{i\star}$, so repeated specimens contribute in proportion to retained cell counts. This all-cell readout differs from $\bar Y_s$ restricted to fibroblasts and from $\widehat w_s$. For healing set $H$ and non-healing set $N$,
\begin{equation}
\widehat\Delta=\frac1{|H|}\sum_{p\in H}S_p-\frac1{|N|}\sum_{p\in N}S_p.\label{eq:p1delta}
\end{equation}
Equation~\eqref{eq:p1delta} is the primary patient-unit contrast; it is separate from the sample-level fibroblast ratio and from the threshold fraction. All $\binom{11}{7}=330$ patient-label allocations are enumerated. To reproduce the recorded result, the implementation retains an add-one correction even for exhaustive enumeration:
\begin{equation}
p_{+}=\frac{1+\sum_{\pi\in\Pi}\ind\{|\widehat\Delta_\pi|\geq|\widehat\Delta|\}}{1+|\Pi|}.\label{eq:p1permutation}
\end{equation}
Equation~\eqref{eq:p1permutation} is the recorded add-one-corrected two-sided permutation probability. It is slightly conservative relative to the uncorrected exhaustive fraction; the add-one convention is usually motivated for randomly drawn permutations, whereas all allocations are enumerated here. Exchangeability is an assumption for an observational label comparison, not a randomization supplied by study design. A 30,000-draw bootstrap resamples patients within outcome arms. Small group sizes and conditional resampling explain why bootstrap intervals and permutation tests need not be dual.

An independent enumeration also reports the exhaustive probability without the extra correction:
\begin{equation}
p_{\rm exact}=\frac{\#\{\pi\in\Pi:|\widehat\Delta_\pi|\geq|\widehat\Delta|\}}{|\Pi|}=\frac{90}{330}=0.272727.\label{eq:p1exact}
\end{equation}
Equation~\eqref{eq:p1exact} includes the observed allocation and numerical ties. It distinguishes exhaustive enumeration from randomly sampled permutations; the recorded add-one comparison $91/331=0.274924$ is retained alongside it. Neither convention makes observational healing labels randomized or tests a causal effect. Table~\ref{tab:p1exact} reports the independent arithmetic check, including the Welch quantities below.

For a symmetric equivalence margin $M$, the Welch two-one-sided-test procedure uses $SE=(s_H^2/n_H+s_N^2/n_N)^{1/2}$ and Welch degrees of freedom $\nu$:
\begin{equation}
p_L=1-F_{t_\nu}\!\left(\frac{\widehat\Delta+M}{SE}\right),\qquad
p_U=F_{t_\nu}\!\left(\frac{\widehat\Delta-M}{SE}\right).\label{eq:p1tost}
\end{equation}
Equation~\eqref{eq:p1tost} gives the two one-sided tests used for a prespecified margin. Both must be below 0.05 to reject non-equivalence. Writing $v_H=s_H^2/n_H$ and $v_N=s_N^2/n_N$ makes the degrees of freedom and margin boundary explicit:
\begin{equation}
\nu=\frac{(v_H+v_N)^2}{v_H^2/(n_H-1)+v_N^2/(n_N-1)},\qquad
M_*=|\widehat\Delta|+t_{\nu,0.95}SE.\label{eq:p1margin}
\end{equation}
Equation~\eqref{eq:p1margin} is the infimum margin containing the two-sided 90\% Welch interval; strict rejection requires $M>M_*$. Independent calculation gives $SE=0.019446$, $\nu=8.113833$ and $M_*=0.065375$ for the 11-patient comparison. The deterministic specimen-unit boundary is 0.079644. A margin selected after seeing these data is a sensitivity summary, not a clinically justified equivalence claim. No clinical minimum important difference has been supplied.

\subsection{Discrimination, anatomy and design simulations}
For held-out prediction scores $r$, AUC is the empirical probability that a healing score exceeds a non-healing score, with half credit for ties:
\begin{equation}
\widehat{\rm AUC}=\frac1{n_Hn_N}\sum_{h\in H,n\in N}\left[\ind\{r_h>r_n\}+\tfrac12\ind\{r_h=r_n\}\right].\label{eq:p1auc}
\end{equation}
Equation~\eqref{eq:p1auc} is the pairwise empirical AUC used for the held-out scores. Topic orientation, module preprocessing and the choice of PCA/NMF components are determined inside each training fold. The topic encoder itself was trained without outcomes on the discovery cohort, including held-out patients' expression: this is an internal discrimination audit of a frozen discovery representation, not fully external representation validation. Bootstrap intervals resample fixed out-of-fold scores with common arm-stratified indices across methods; the entire fitting/orientation process is rerun for label permutations. All allocations are used for topic and module readouts; PCA and NMF use 200 sampled relabellings. Sample-level comparisons use 14 specimens and patient-level comparisons use 11 mapped patients; they are not interchangeable.

PCA scaling and basis fitting, and NMF fitting, use training rows only. The historical pooled out-of-fold AUC nevertheless compares scores from different fitted folds; those scores are not calibrated clinical risks. Separate PCA and NMF fits can change the selected coordinate, origin and scale between folds. AUC is invariant to a common increasing transformation, but not to separate fold-specific transformations. These raw-coordinate AUCs and their fixed-score bootstrap contrasts are retained as historical scoring diagnostics, not evidence for a representation ranking.

The revised patient discrimination audit instead leaves out one healing patient $h$ and one non-healing patient $n$ together. Each of the $7\times4=28$ folds fits one scoring function $r^{(-h,-n)}$ to the remaining nine patients (six healing, three non-healing), and that unchanged function scores both excluded patients:
\begin{equation}
\widehat A_{\rm pair}=\frac1{28}\sum_{h\in H,n\in N}\left[\ind\{r^{(-h,-n)}_h>r^{(-h,-n)}_n\}+\tfrac12\ind\{r^{(-h,-n)}_h=r^{(-h,-n)}_n\}\right].\label{eq:p1pairauc}
\end{equation}
Equation~\eqref{eq:p1pairauc} averages within-model pair orderings; no raw scores from different fitted models are compared. Patient expression inputs are cell-count-weighted means of per-cell proportions on the fixed panel. Topic orientation, module gene standardization and orientation, PCA standardization and fitting, and NMF fitting and component selection use training patients only. The PCA rule selects the largest absolute training-arm mean difference among five components. NMF uses five components and tolerance $10^{-4}$; the initial 1,000-iteration cap was raised uniformly to 10,000 after three training fits reached the cap, without using held-pair discrimination to select it. The initial run is retained, and no final fit emitted a convergence warning. Fitted basis rows are normalized to unit Euclidean norm before independent nonnegative-least-squares projection of each row. The component is selected by the absolute training-arm mean difference divided by the training-patient SD, making selection invariant to positive component rescaling. The direction is set by the training contrast, and exact ties receive half credit. The fixed gene panel, topic representation and study-specific analysis choices remain discovery-derived. The 28 pairs share 11 patients and overlapping training sets, so they are not 28 independent outcomes; no pair-binomial uncertainty, fixed-score bootstrap interval or significance probability is supplied for this revised estimand. Changing component normalization and selection as well as fold structure means the difference from the historical AUC cannot identify the contribution of one correction alone. Reusing the same cohort to inspect readouts limits confirmatory interpretation.

The paired anatomical control tests foot-minus-forearm differences in 10 mapped patients using all $2^{10}$ sign allocations and the same add-one convention. Across 15 topics, Benjamini--Hochberg adjustment is $q_{(i)}=\min\{1,\min_{j\geq i}15p_{(j)}/j\}$, with the usual assumptions required for false-discovery-rate interpretation.

The historical high-state-fraction power calculation uses 9 versus 5 specimens, Gaussian location shifts with pooled SD 0.283 and a two-sided Mann--Whitney test. Its 24.9\% power and $\operatorname{MDE}_{0.8}=0.519$ apply to this approximate design, not to the 11-patient all-cell contrast. Gaussian simulation can leave the unit interval, so it is a design-scale sensitivity approximation. A separate next-cohort simulation uses specimen-derived arm SDs of deterministic all-cell mean scores, 0.05170 and 0.01686, Welch testing or Welch equivalence testing, and 4,000 draws per candidate size. These descriptive spreads include repeated specimens and are not independent-patient variance estimates. The candidate totals are approximate conditional calculations, not validated recruitment targets (Table~\ref{tab:p1design}).

\subsection{External projections and limits of identification}
Every external cohort is projected into the frozen panel/model without outcome-guided retuning. The immune and B/plasma gates exclude promiscuous markers such as CD74, CD37, CD38 and SDC1 and assign non-overlapping compartments. Immune fraction uses all retained cells as its denominator; B/plasma share within immune cells uses only that compartment. Their correlations estimate different quantities. Negative depth correlations cannot by themselves exclude depth confounding, because a marginal sign is not a conditional adjustment.

GSE231643 supplies eight title-labelled specimens but lacks an authoritative patient mapping. GSE241132 contributes an acute-wound, three-donor construct summary. GSE223964, GSE245703 and GSE268834 lack qualifying patient-level healing endpoints, and GSE248247 has only two samples. We do not pool these correlations or describe them as independent prognostic validation. No analysis identifies a treatment effect, a time-to-healing model or a causal effect of immune infiltration.

\subsection{Observed-cell and expression displays}
The biological displays reuse the saved 15-topic cell proportions and operational lineage assignments. For visualization only, each proportion vector is transformed to centered log ratios,
\begin{equation}
c_{ik}=\log\theta_{ik}-K^{-1}\sum_{\ell=0}^{K-1}\log\theta_{i\ell}.\label{eq:p1display}
\end{equation}
Equation~\eqref{eq:p1display} removes the geometric mean of the positive coordinates before a Euclidean-distance UMAP with 30 neighbors, minimum distance 0.3 and random seed 17. All 76,461 cells enter the embedding; a uniform subset of 2,000 cells is drawn for the printed vector panels, using the same cells and coordinates for lineage and loading colors. The sampling affects only display. Neither embedding separation nor displayed point density enters the outcome tests, gate, threshold or model fitting. Two-dimensional maps can distort relationships in the input space and do not establish cell types or trajectories.

Raw discovery matrices are reloaded under the original tissue, disease and cell-quality criteria. We require the retained sample and arm order and all recomputed marker-based labels to match the saved cell table. For each displayed gene, dot size is the fraction of cells with a positive count; color is the mean of cellwise $\log(1+10^4 x_{ig}/\sum_h x_{ih})$, scaled from its smallest to largest lineage mean. The denominator uses the source gene matrix, not only the frozen modeling panel. The 18-gene display includes compartment markers and programme-associated genes without an outcome-based selection step. It illustrates measured expression and cannot independently validate labels derived from some of the same markers. Specimen composition retains each specimen as a separate column, and loading histograms use all assigned fibroblasts in fixed bins of width 0.02.

\subsection{Patient-unit extension of the lineage readouts}
The retrospective extension fixes three readouts before computing their new outcome comparisons: within-fibroblast mean loading, the fraction above the existing mixture midpoint, and fibroblast composition among all retained cells. Let $\mathcal S_p$ denote patient $p$'s specimens, $n_s^F$ their fibroblast counts and $n_s^A$ their all-cell counts. With specimen means $m_s$, threshold fractions $w_s$ and composition fractions $f_s=n_s^F/n_s^A$, the pooled patient summaries are
\begin{equation}
M_p=\frac{\sum_{s\in\mathcal S_p}n_s^F m_s}{\sum_{s\in\mathcal S_p}n_s^F},\quad
W_p=\frac{\sum_{s\in\mathcal S_p}n_s^F w_s}{\sum_{s\in\mathcal S_p}n_s^F},\quad
F_p=\frac{\sum_{s\in\mathcal S_p}n_s^A f_s}{\sum_{s\in\mathcal S_p}n_s^A}.\label{eq:p1lineagepatients}
\end{equation}
Equation~\eqref{eq:p1lineagepatients} pools the appropriate denominators within patient while giving each of the 11 outcome patients equal weight in the group contrast. The representation, annotations and loading cut are unchanged. We enumerate all 330 allocations of seven healed labels among 11 patient rows, counting absolute contrast ties within a numerical tolerance of $10^{-12}$. For this new exhaustive enumeration the probability is the exceeding count divided by 330, without the add-one convention retained in the historical all-cell analysis. The three new probabilities form one Benjamini--Hochberg family. Observational group exchangeability is an assumption, not random assignment of healing outcomes.

Percentile intervals use 10,000 patient resamples separately within the seven healed and four non-healed groups, with seed 17. They are small-sample descriptive intervals, not externally calibrated prediction intervals. Sensitivities replace within-patient count weights by equal specimen weights, omit each patient in turn, and require every retained specimen to have at least 100 fibroblasts. Common random draws are used when the patient arrays coincide. We report when an eligibility sensitivity removes no specimen. The 20-patient descriptive table includes non-diabetic participants, but only the 11 DFU patients enter healing contrasts; no healthy participant is assigned a healing outcome.

\subsection{Patient distributions, cell yield and repeated sampling}
To assess the loading distribution without selecting one state threshold, let $F_p(c)$ be the empirical cumulative distribution of all recovered fibroblast loadings pooled within patient $p$. Each patient contributes a normalized distribution regardless of cell yield. We compare the two equally weighted patient-group distributions using
\begin{equation}
T=\sup_{c\in\mathcal C}\left|\frac1{7}\sum_{p:H_p=1}F_p(c)-\frac1{4}\sum_{p:H_p=0}F_p(c)\right|,\label{eq:p1ecdf}
\end{equation}
where $\mathcal C$ contains every distinct observed loading. Equation~\eqref{eq:p1ecdf} uses the same 11 patients and frozen scores. Its exact probability enumerates all 330 patient-label allocations; this is a patient-level permutation statistic, not a cell-level Kolmogorov--Smirnov test. We also report high-fraction contrasts at cuts 0.05, 0.10, 0.15, 0.20, 0.30, 0.40, 0.50 and 0.60, together with the original midpoint. Comparing each absolute contrast with the permutation distribution of $T$ gives a max-statistic probability across the complete loading support. These retrospective sensitivities and their unadjusted probabilities are shown in Table~\ref{tab:p1thresholdpatient}.

Cell-yield sensitivity uses 250 draws of exactly 100 fibroblasts without replacement from each of the 14 DFU specimens, with random seed 29. Each draw averages specimen means and high-state fractions equally within patient before recomputing the 7-versus-4 contrast and exact probability. This changes the contribution of recovered cells, without changing clinical labels or the frozen threshold. Fibroblast composition is not recomputed from this deliberately selected cell subset. Separately, three patients have two specimens each; all $2^3=8$ choices of one specimen per patient are evaluated for all three original readouts. Table~\ref{tab:p1samplingsensitivity} gives complete contrast and probability ranges. The draws and specimen choices assess finite-cell and sampling sensitivity, not new patients or confidence bounds.


Repeated-specimen diagnostics retain all five mapped pairs, including three DFU pairs. For each pair we compare fibroblast means, high-state fractions and the supremum difference between empirical loading distributions. Within each specimen, 250 seeded random partitions create disjoint fibroblast halves without replacement; their absolute mean differences provide a conditional cell-sampling reference. We compare each observed between-specimen difference with both within-specimen 97.5th percentiles. These partitions are not biological replicates, patient confidence intervals or an intraclass correlation estimate (Table~\ref{tab:p1repeats}).


\section*{Data and reproducibility}
The discovery series is GSE165816. External projections use GSE231643, GSE241132, GSE223964, GSE245703, GSE268834 and GSE248247; each accession is linked in Table~\ref{tab:p1cohorts}. This is a secondary analysis of publicly released data, with no new recruitment or intervention. The original studies govern participant consent and source-study ethical approvals. Shared series-level laboratory metadata do not identify sample-specific sequencing batches. The mathematical definitions, analysis units, numerical summaries and figure source data specify the present analysis; no unavailable patient identifiers, healing times or independent clinical validation are imputed.
Analysis code is released under the MIT License. Manuscripts, figures and source-study data remain under their respective rights and repository terms. The accompanying software archive, authored by Zeyu Fu, is available at \url{https://github.com/PeterPonyu/wound-measurement-construct}. An independently archived software release is available at \url{https://doi.org/10.5281/zenodo.22875656}.
\FloatBarrier
\bibliographystyle{abbrvnat}
\bibliography{references}
\supplement
Table~\ref{tab:p1cohorts} fixes the role of each accession. Specimen counts, donor counts and cell counts answer different questions. The 54-sample source series is not the 25-sample frozen foot-skin subset, and neither is the 11-patient outcome subset. Cells in an external specimen are not independent replications of that specimen's clinical label.
\begin{table}[!htbp]\centering\small
\caption{Cohort roles, analysis units and interpretation limits. GEO accessions identify the public source studies. Counts refer to specimens, mapped patients or donors as explicitly stated, rather than independent cells. DFU denotes diabetic foot ulcer and QC denotes quality control. The original discovery series, selected foot-skin representation and outcome subset have different denominators; external projections test construct associations or feasibility without a qualifying patient-level healing validation.}\label{tab:p1cohorts}
\begin{tabularx}{\linewidth}{@{}p{26mm}p{32mm}L@{}}\toprule
Cohort & Units & Role and limit \\\midrule
\href{https://www.ncbi.nlm.nih.gov/geo/query/acc.cgi?acc=GSE165816}{GSE165816} & 54 original specimens; 25 selected foot specimens & Discovery representation; 14 DFU specimens map to 11 outcome patients. The original series maps to 27 participants. \\
\href{https://www.ncbi.nlm.nih.gov/geo/query/acc.cgi?acc=GSE231643}{GSE231643} & 8 specimens & Title-derived healing labels (5, 3); no authoritative independent patient mapping. \\
\href{https://www.ncbi.nlm.nih.gov/geo/query/acc.cgi?acc=GSE241132}{GSE241132} & 3 donors; 12 timepoint samples & Healthy acute-wound construct audit; no DFU prognosis inference. \\
\href{https://www.ncbi.nlm.nih.gov/geo/query/acc.cgi?acc=GSE223964}{GSE223964} & 8 specimens & Frozen projection and exclusive immune-compartment audit; no qualifying healing outcome. \\
\href{https://www.ncbi.nlm.nih.gov/geo/query/acc.cgi?acc=GSE245703}{GSE245703} & 12 specimens & Frozen projection and immune-compartment audit; no qualifying healing outcome. \\
\href{https://www.ncbi.nlm.nih.gov/geo/query/acc.cgi?acc=GSE268834}{GSE268834} & 8 specimens & Descriptive diabetes-group and immune-context contrasts, not healing validation. \\
\href{https://www.ncbi.nlm.nih.gov/geo/query/acc.cgi?acc=GSE248247}{GSE248247} & 2 specimens; 358 QC cells & Low-yield projection check from 561 raw cells; correlations not estimated. \\\bottomrule
\end{tabularx}\end{table}

Discovery lineage scores average gene-wise standardized $\log(1+10^4x_{ig}/\sum_hx_{ih})$ values within marker sets; the highest lineage score assigns the label. The fibroblast markers are COL1A1, COL1A2, PDGFRA, TWIST2, COL6A3, MMP2 and CCDC80. The circularity check concerns overlap with the tested topic's specified top-gene list; it does not imply independence from all genes contributing to the neural representation. The loading is a learned function of the full panel.

In external audits, immune assignment is PTPRC detection or detection of at least two markers within a T, NK, myeloid, B or plasma panel. These panels respectively contain CD3D, CD3E, CD2, TRBC1, TRBC2, NKG7, GNLY; LYZ, CD68, CD14, ITGAX, FCER1G, TYROBP, AIF1; MS4A1, CD79A, CD79B, BANK1; and MZB1, JCHAIN, DERL3, TNFRSF17. Fibroblast-like cells require at least two of COL1A1, COL1A2, COL3A1, DCN, LUM, PDGFRA, THY1 and no immune assignment. B/plasma cells are immune-assigned cells detecting at least one B or plasma marker. All remaining cells form the other compartment. Unlike the discovery lineage list, this external structural list includes topic-associated genes; it is a descriptive operational gate and is not evidence of gene-disjoint external validation.

The disjoint immune, fibroblast-like and other compartments sum to one. B/plasma is a subset of immune, so its within-immune share must not be added as a fourth exclusive compartment. The older overlapping marker-detection fractions are not used for the current external associations.

For distributional descriptions, pure high requires the sample's fibroblast 10th percentile to be at least 0.10; pure low requires its 90th percentile to be at most 0.15; mixed requires a 10th percentile at most 0.05 and a 90th percentile at least 0.40. Remaining samples are labelled other. These operational categories do not prove distinct biological attractors.

Let the reference and candidate $k$-nearest-neighbour graphs have component labels $r_i,c_i$. If $n_{ab}$ counts cells jointly in reference component $a$ and candidate component $b$, and $n_a=\sum_bn_{ab}$, the implementation computes reachability preservation exactly as
\begin{equation}
R_{\rm reach}=\frac{\sum_{a,b}n_{ab}^2-N}{\sum_an_a^2-N},\label{eq:p1reach}
\end{equation}
Equation~\eqref{eq:p1reach} is assigned value one when the denominator is zero. This sparse component-count calculation avoids materializing all pairwise reachability values. It is a connectivity statistic: two connected graphs can give one despite severely distorted local geometry. The separate edge-survival score is reported but is not included in the binary integrity condition.

Mean cell perplexity is $\overline P=N^{-1}\sum_i\exp[-\sum_k\theta_{ik}\log\theta_{ik}]$. The recorded integrity condition is $R_{\rm reach}>0.70$ and $\overline P\leq0.95K$. A fixed 6,000-cell raw-count subset was permuted gene by gene across cells. This preserves each gene\textquotesingle{}s observed value multiset but not row-wise library sizes; it raises $\overline P$ from 3.44 of 15 to 11.23 of 15 and lowers the mean HHI from 0.574 to 0.147. The shuffled input is therefore less concentrated and more structureless under this coordinate summary, but the experiment is not a complete technical null because library-size and detection-rate variation also narrow. It does not validate every topological feature or every cut at the fixed reachability bound. The $K$ and state-cut grids in the main text do not constitute a sensitivity scan of the 0.70 connectivity bound.

Every interval names its resampling unit. The primary semantic bootstrap resamples the 20 mapped patients after fibroblast-count-weighted pooling; the earlier sample-row bootstrap remains a historical diagnostic, not the primary uncertainty estimate. The all-cell outcome bootstrap resamples patients within arms after replicate collapse. Historical pooled-fold AUC intervals resample fixed out-of-fold predictions; they do not refit the encoder or reproduce the entire learning algorithm on each bootstrap draw, and do not apply to the revised same-model held-pair estimand. Permutation inference reruns outcome-dependent score orientation and representation selection, which is why the appropriate null AUC need not be 0.5. The unadjusted $p=0.0424$ for the sample-level module is exploratory across several readouts and is not a confirmatory multiplicity-controlled endpoint.

The next-study design calculation uses sample-derived arm spreads. This creates uncertainty not reflected by its Monte Carlo standard error; at estimated power 0.8 and 4,000 simulations, the conditional standard error is about $\sqrt{0.8(0.2)/4000}=0.0063$. Candidate equivalence designs assume a true difference of zero. Bounded topic scores are approximated by Gaussian locations and variances, and proposed totals near 80\% power are approximate search outputs rather than guaranteed minima.

For the exploratory AUC-design simulator, a unit-variance binormal score with class-mean separation $d$ has $\mathrm{AUC}=\Phi(d/\sqrt2)$, hence $d=\sqrt2\Phi^{-1}(\mathrm{AUC})$. Correlation is imposed on within-class noise, not on the full mean-shifted scores. Its standard-error approximation is estimated across simulated AUC differences; it is a model-based planning exercise, not a validated patient-level test. Neither these AUC calculations nor an intrapatient design-effect formula licenses treating repeated biopsies as independent prognosis cases.

Tables~\ref{tab:p1patients}--\ref{tab:p1external} present the mapped patient scores, discrimination intervals, cross-method associations, conditional design calculations, semantic uncertainty and external immune-context summaries. Their denominators and limits accompany the values so that the three distinct readouts---all-cell mean, fibroblast mean and high-state fraction---remain identifiable.
Tables~\ref{tab:p1covariates} and~\ref{tab:p1ambient} retain the complete covariate screen and immune/placebo residual summaries. Tables~\ref{tab:p1lineage}, \ref{tab:p1aucdifferences}, \ref{tab:p1anatomy} and~\ref{tab:p1equivalence} distinguish the lineage, discrimination, paired-anatomy and equivalence estimands. Tables~\ref{tab:p1cuts} and~\ref{tab:p1capacity} retain every tested threshold and topic capacity, including adverse sensitivities. Table~\ref{tab:p1aucdesign} records the exploratory AUC-gap simulation. The separate arithmetic conversion from specimens to patients under assumed intrapatient correlation is not a validated patient-prognosis design and is not used as a recruitment recommendation.
{\renewcommand{\arraystretch}{1.00}{\small
\begin{longtable}{@{}p{22mm}p{29mm}p{39mm}p{26mm}p{49mm}@{}}
\caption{Author-mapped GSE165816 DFU patient scores (7 healed, 4 not healed). Each score is the mean frozen fibroblast-associated topic loading over all retained cells in the listed foot specimens; pooling cells weights repeated specimens by retained cell count. Cells are counts, not independent outcome observations. Patient identifiers are study-derived pseudonyms and G labels identify specimens. The outcome is the source-study category within 12 weeks, not an individual healing day. Scores are on the topic-loading scale from zero to one.}\label{tab:p1patients}\\
\toprule
Patient & Outcome & Specimens & Cells & All-cell mean \\
\midrule\endfirsthead
\multicolumn{5}{l}{\tablename\ \thetable\ (continued)}\\
\toprule
Patient & Outcome & Specimens & Cells & All-cell mean \\
\midrule\endhead
\bottomrule\endfoot
S01 & not healed & G6 & 2279 & 0.0498 \\*
S02 & healed & G23, G4 & 7495 & 0.1496 \\
S03 & healed & G2 & 3951 & 0.0185 \\
S04 & healed & G15 & 2250 & 0.0758 \\
S05 & not healed & G9 & 4290 & 0.0346 \\
S06 & not healed & G33, G34 & 7681 & 0.0103 \\
S07 & not healed & G39 & 2481 & 0.0277 \\
S08 & healed & G7, G8 & 4192 & 0.0754 \\
S09 & healed & G49 & 2928 & 0.0534 \\
S10 & healed & G42 & 1838 & 0.0167 \\*
S11 & healed & G45 & 3374 & 0.0297 \\
\end{longtable}
}

{\small
\begin{longtable}{@{}p{25mm}p{20mm}p{49mm}p{37mm}p{34mm}@{}}
\caption{Historical scoring diagnostics across 9 healing and 5 non-healing GSE165816 specimens. AUC is the area under the ROC curve for pooled leave-one-sample-out scores; higher values favour healing. The percentile 95\% intervals use 4,000 arm-stratified resamples of fixed out-of-fold scores. One-sided permutation probabilities and null means rerun fold-dependent orientation and fitting. PCA is principal component analysis; NMF is nonnegative matrix factorization. Their raw coordinates are not calibrated across fitted folds, so these estimates cannot rank representations. The frozen topic encoder used discovery expression, and repeated patients limit independence. None of the six historical paired AUC-difference intervals excludes zero.}\label{tab:p1sampleauc}\\
\toprule
Readout & AUC & Bootstrap 95\% interval & Permutation p & Null mean \\
\midrule\endfirsthead
\multicolumn{5}{l}{\tablename\ \thetable\ (continued)}\\
\toprule
Readout & AUC & Bootstrap 95\% interval & Permutation p & Null mean \\
\midrule\endhead
\bottomrule\endfoot
Fibroblast topic & 0.822 & 0.556 to 1.000 & 0.0604 & 0.562 \\*
Module & 0.844 & 0.578 to 1.000 & 0.0424 & 0.497 \\
PCA & 0.533 & 0.156 to 0.889 & 0.4527 & 0.480 \\*
NMF & 0.878 & 0.655 to 1.000 & 0.0597 & 0.553 \\
\end{longtable}
}

{\small
\begin{longtable}{@{}p{45mm}p{25mm}p{59mm}p{38mm}@{}}
\caption{Historical patient-unit scoring diagnostics and paired anatomical sensitivity within GSE165816. The first four rows are pooled leave-one-patient-out AUCs for 7 healing and 4 non-healing patients, with 4,000 arm-stratified bootstrap draws of fixed scores and one-sided label-permutation probabilities. PCA and NMF denote principal component analysis and nonnegative matrix factorization; their coordinates are not comparable across fitted folds. These historical estimates are not uncertainty summaries for the revised held-pair procedure and do not rank representations. The final row is the foot-minus-forearm all-cell fibroblast-associated topic mean difference in 10 paired patients, with a 4,000-draw paired bootstrap interval and a two-sided exhaustive sign-flip probability using add-one correction. Its Benjamini--Hochberg q value across 15 topics is 0.0402. These row types have different estimands; neither supplies external validation.}\label{tab:p1patientauc}\\
\toprule
Readout & Estimate & Bootstrap 95\% interval & Permutation p \\
\midrule\endfirsthead
\multicolumn{4}{l}{\tablename\ \thetable\ (continued)}\\
\toprule
Readout & Estimate & Bootstrap 95\% interval & Permutation p \\
\midrule\endhead
\bottomrule\endfoot
Fibroblast topic & 0.7143 & 0.3571 to 1.0000 & 0.2991 \\*
Module & 0.7857 & 0.5000 to 1.0000 & 0.1662 \\
PCA & 0.5000 & 0.0000 to 1.0000 & 0.5025 \\
NMF & 0.1250 & 0.0000 to 0.4286 & 0.9751 \\*
foot minus forearm & 0.0476 & 0.0119 to 0.0916 & 0.0107 \\
\end{longtable}
}

\Needspace{155pt}
{\small
\begin{longtable}{@{}p{53mm}p{53mm}p{65mm}@{}}
\caption{Spearman correlations between the four representative sample-level prediction scores across the same 14 DFU specimens (9 healing, 5 non-healing). Fibroblast topic is the frozen all-cell loading, Module the specified fibro-inflammatory score, PCA principal component analysis and NMF nonnegative matrix factorization. Each row is one pair of readouts; values are descriptive point estimates without confidence intervals or multiplicity-adjusted tests. Related scores do not establish equivalent discrimination or rank methods.}\label{tab:p1correlations}\\
\toprule
Readout A & Readout B & Spearman correlation \\
\midrule\endfirsthead
\multicolumn{3}{l}{\tablename\ \thetable\ (continued)}\\
\toprule
Readout A & Readout B & Spearman correlation \\
\midrule\endhead
\bottomrule\endfoot
Fibroblast topic & Module & 0.7099 \\*
Fibroblast topic & PCA & 0.1736 \\
Fibroblast topic & NMF & 0.6067 \\
Module & PCA & 0.4022 \\
Module & NMF & 0.6756 \\*
PCA & NMF & 0.3200 \\
\end{longtable}
}

{\small
\begin{longtable}{@{}p{50mm}p{30mm}p{20mm}p{30mm}p{35mm}@{}}
\caption{Conditional next-cohort simulations for all-cell fibroblast-associated topic mean loading. Rows assume hypothetical independent observations, no missing outcomes, healing allocation near 9/14, and arm SDs 0.05170 and 0.01686 estimated from deterministic discovery specimens. These descriptive spreads include repeated specimens and are not patient-population variance estimates. Two-sided Welch tests assess candidate mean differences; Welch two-one-sided tests assess equivalence under a true difference of zero, at alpha=0.05. Total and arm columns give the simulated observation counts; power uses 4,000 draws per candidate size. Values near 0.80 have Monte Carlo error and are approximate search outputs. Candidate differences and margins are inputs, not clinical minimum important differences or validated recruitment targets.}\label{tab:p1design}\\
\toprule
Design & Difference / margin & Total & Healing / other & Estimated power \\
\midrule\endfirsthead
\multicolumn{5}{l}{\tablename\ \thetable\ (continued)}\\
\toprule
Design & Difference / margin & Total & Healing / other & Estimated power \\
\midrule\endhead
\bottomrule\endfoot
Detect difference & 0.02 & 100 & 64 / 36 & 0.797 \\*
Detect difference & 0.03 & 47 & 30 / 17 & 0.821 \\
Detect difference & 0.0457 & 21 & 14 / 7 & 0.811 \\
Detect difference & 0.06 & 14 & 9 / 5 & 0.804 \\
Detect difference & 0.08 & 9 & 6 / 3 & 0.807 \\
Detect difference & 0.1 & 8 & 5 / 3 & 0.880 \\
Equivalence under zero difference & 0.02 & 110 & 71 / 39 & 0.815 \\
Equivalence under zero difference & 0.03 & 49 & 32 / 17 & 0.826 \\
Equivalence under zero difference & 0.05 & 20 & 13 / 7 & 0.830 \\
Equivalence under zero difference & 0.075 & 11 & 7 / 4 & 0.880 \\*
Equivalence under zero difference & 0.1 & 8 & 5 / 3 & 0.915 \\
\end{longtable}
}

{\small
\begin{longtable}{@{}p{44mm}p{38mm}p{51mm}p{34mm}@{}}
\caption{Exploratory binormal planning simulation for paired AUC differences. Two candidate AUCs are centred on 0.80, separated by the stated gap, with within-class score-noise correlation 0.5045 and healing allocation near 9/14. Each searched total uses 4,000 simulation draws and a model-based normal approximation for the paired difference at alpha=0.05. Totals assume one observation per independent patient and no missing outcome. The gaps are hypothetical inputs; the approximate 80\% power search does not establish actual clinical performance, calibrated patient-test power or a validated recruitment target.}\label{tab:p1aucdesign}\\
\toprule
Candidate AUC gap & Total observations & Healing / non-healing & Estimated power \\
\midrule\endfirsthead
\multicolumn{4}{l}{\tablename\ \thetable\ (continued)}\\
\toprule
Candidate AUC gap & Total observations & Healing / non-healing & Estimated power \\
\midrule\endhead
\bottomrule\endfoot
0.05 & 706 & 454 / 252 & 0.800 \\*
0.1 & 181 & 116 / 65 & 0.814 \\
0.15 & 85 & 55 / 30 & 0.797 \\*
0.2 & 50 & 32 / 18 & 0.800 \\
\end{longtable}
}

{\small
\begin{longtable}{@{}p{63mm}p{28mm}p{39mm}p{37mm}@{}}
\caption{Association between mean fibroblast-associated topic loading within fibroblasts and the threshold-defined high-state fraction across 25 GSE165816 specimens. The Spearman row gives the observed correlation and percentile 95\% limits from 10,000 sample-row bootstrap draws. The leave-one-out row gives the minimum and maximum of 25 recalculations, not a confidence interval; a dash means no single point estimate applies. A separate 10,000-draw permutation analysis gave a two-sided add-one probability of 0.00010 and a central null interval from -0.4025 to 0.4072. Samples may share patients, and the two summaries share cells.}\label{tab:p1semantics}\\
\toprule
Summary & Estimate & Lower endpoint & Upper endpoint \\
\midrule\endfirsthead
\multicolumn{4}{l}{\tablename\ \thetable\ (continued)}\\
\toprule
Summary & Estimate & Lower endpoint & Upper endpoint \\
\midrule\endhead
\bottomrule\endfoot
spearman rho & 0.9325 & 0.7887 & 0.9828 \\*
leave one out rho & — & 0.9235 & 0.9566 \\
\end{longtable}
}

{\small
\begin{longtable}{@{}p{29mm}p{18mm}p{29mm}p{29mm}p{29mm}p{29mm}@{}}
\caption{External sample-level Spearman correlations with the frozen fibroblast-associated topic mean under exclusive lineage gates. Immune fraction divides immune-assigned cells by all retained cells; B/plasma share divides B/plasma cells by immune-assigned cells. Detected genes and unique molecular identifiers (UMIs) are separate per-sample medians over QC cells. All entries are descriptive point estimates without confidence intervals. Specimens are not verified independent patients, and no listed cohort supplies a qualifying patient-level healing outcome.}\label{tab:p1external}\\
\toprule
Cohort & Samples & Immune fraction & B/plasma share & Detected genes & UMIs \\
\midrule\endfirsthead
\multicolumn{6}{l}{\tablename\ \thetable\ (continued)}\\
\toprule
Cohort & Samples & Immune fraction & B/plasma share & Detected genes & UMIs \\
\midrule\endhead
\bottomrule\endfoot
GSE223964 & 8 & -0.024 & -0.071 & -0.429 & -0.476 \\*
GSE245703 & 12 & 0.846 & 0.077 & -0.427 & -0.308 \\*
GSE268834 & 8 & 0.714 & -0.500 & 0.238 & 0.262 \\
\end{longtable}
}

{\small
\begin{longtable}{@{}p{76mm}p{20mm}p{25mm}p{22mm}p{22mm}@{}}
\caption{Complete exploratory covariate screen in 25 discovery specimens. Each row correlates the fibroblast high-state fraction, using the historical rule loading at least 0.184, with one of 31 covariates. Lineage fractions use all retained cells; technical summaries and topic means use fibroblasts. Values are Spearman rho, two-sided asymptotic p and Benjamini--Hochberg q across all 31 tests, ordered by absolute rho; 14 rows have q below 0.05. UMI means unique molecular identifier. The G label is an identifier, not an established batch or time variable. Repeated patients invalidate the naive independent-row premise of these historical p values; their BH transforms do not establish false-discovery-rate control or 14 independent discoveries. Shared cells and compositional constraints further limit interpretation; these are descriptive row-level summaries, not causal or healing-validation tests.}\label{tab:p1covariates}\\
\toprule
Covariate & Specimens & Spearman rho & Raw p & BH q \\
\midrule\endfirsthead
\multicolumn{5}{l}{\tablename\ \thetable\ (continued)}\\
\toprule
Covariate & Specimens & Spearman rho & Raw p & BH q \\
\midrule\endhead
\bottomrule\endfoot
B/plasma fraction & 25 & 0.8279 & 3.26e-07 & 1.01e-05 \\*
Fibroblast CFD/APOD mean & 25 & -0.7664 & 7.95e-06 & 0.000123 \\
Fibroblast PTPRC-positive fraction & 25 & 0.7267 & 3.88e-05 & 0.000401 \\
Fibroblast median UMI count & 25 & 0.6632 & 0.000302 & 0.00234 \\
Fibroblast IGFBP7/TAGLN mean & 25 & 0.6514 & 0.000421 & 0.00261 \\
Sweat-gland fraction & 25 & -0.5832 & 0.00222 & 0.0114 \\
Fibroblast FOS/IGFBP7 mean & 25 & -0.5687 & 0.00302 & 0.0134 \\
Fibroblast DCN/CFD mean & 25 & -0.5592 & 0.00366 & 0.0142 \\
Fibroblast IGFBP7/IGFBP5 mean & 25 & 0.5316 & 0.00624 & 0.0215 \\
Fibroblast COL1A1/COL1A2 mean & 25 & 0.5238 & 0.00721 & 0.0223 \\
Keratinocyte fraction & 25 & -0.4836 & 0.0143 & 0.0404 \\
Fibroblast HLA-DRA/CD74 mean & 25 & 0.4726 & 0.0171 & 0.0441 \\
T-cell fraction & 25 & 0.4647 & 0.0193 & 0.0459 \\
Fibroblast TAGLN/ACTA2 mean & 25 & 0.4552 & 0.0222 & 0.0492 \\
Melanocyte fraction & 25 & -0.4442 & 0.0261 & 0.054 \\
Fibroblast JUNB/DNAJB1 mean & 25 & 0.4182 & 0.0375 & 0.0726 \\
Mast-cell fraction & 25 & -0.3977 & 0.0489 & 0.0893 \\
Myeloid fraction & 25 & 0.3891 & 0.0546 & 0.094 \\
Numeric G specimen label & 25 & -0.3607 & 0.0765 & 0.125 \\
Endothelial fraction & 25 & -0.2347 & 0.259 & 0.401 \\
Total retained cells & 25 & -0.2276 & 0.274 & 0.404 \\
Fibroblast DNAJB1/FOS mean & 25 & 0.1977 & 0.344 & 0.484 \\
Fibroblast KRT14/KRT10 mean & 25 & 0.1457 & 0.487 & 0.63 \\
Fibroblast dissociation score & 25 & 0.1449 & 0.489 & 0.63 \\
Lymphatic endothelial fraction & 25 & -0.1331 & 0.526 & 0.63 \\
Fibroblast fraction & 25 & 0.1292 & 0.538 & 0.63 \\
Fibroblast FOS/JUNB mean & 25 & -0.1260 & 0.548 & 0.63 \\
Retained fibroblasts & 25 & 0.0851 & 0.686 & 0.759 \\
Pericyte/smooth-muscle fraction & 25 & -0.0473 & 0.823 & 0.879 \\
Fibroblast IGFBP7/SPARCL1 mean & 25 & -0.0354 & 0.866 & 0.895 \\*
Fibroblast DCD/SAT1 mean & 25 & -0.0236 & 0.911 & 0.911 \\
\end{longtable}
}

{\small
\begin{longtable}{@{}p{64mm}p{13mm}p{22mm}p{23mm}p{23mm}p{20mm}@{}}
\caption{Immune-transcript exclusion and residual-mixture controls. BIC denotes Bayesian information criterion, fitted to sample high-state fractions or their regression residuals; the difference is BIC for one Gaussian minus BIC for two. Positive differences favour two components. The first two rows share 24 matched specimens and have fraction correlation 0.9887. Residual regressions use all 25 specimens: immune covariates are B/plasma, T-cell and myeloid fractions plus immune-transcript positivity; placebos are pericyte/smooth-muscle, endothelial, melanocyte and sweat-gland fractions. R squared is variance explained by each regression and is inapplicable to unadjusted rows (dash). Both residual differences are negative, so the analysis does not establish immune-specific separation.}\label{tab:p1ambient}\\
\toprule
Summary & n & R squared & BIC one & BIC two & Difference \\
\midrule\endfirsthead
\multicolumn{6}{l}{\tablename\ \thetable\ (continued)}\\
\toprule
Summary & n & R squared & BIC one & BIC two & Difference \\
\midrule\endhead
\bottomrule\endfoot
Before immune-transcript exclusion & 24 & — & 12.417 & -28.569 & 40.986 \\*
After immune-transcript exclusion & 24 & — & 5.342 & -51.341 & 56.683 \\
Immune-covariate residuals & 25 & 0.4953 & -3.900 & 0.031 & -3.931 \\*
Placebo-covariate residuals & 25 & 0.3549 & 2.235 & 10.433 & -8.198 \\
\end{longtable}
}

{\small
\begin{longtable}{@{}p{29mm}p{49mm}p{24mm}p{42mm}p{21mm}@{}}
\caption{Historical sample-level lineage and loading contrasts. Ratios divide the healing-arm mean by the non-healing-arm mean; intervals are percentile 95\% limits from 20,000 within-arm sample bootstraps, and p values are two-sided Mann--Whitney tests. GSE165816 uses 9 versus 5 DFU specimens; GSE231643 uses 5 versus 3 title-labelled specimens. The discovery ratios come from the saved lineage-projection analysis and are distinct from the later deterministic all-cell patient contrast. Fibroblast fraction uses all cells as denominator; fibroblast mean uses only assigned fibroblasts. These are separate estimands. Repeated discovery patients and unavailable external patient mappings prevent independent outcome validation.}\label{tab:p1lineage}\\
\toprule
Cohort & Readout & Mean ratio & Bootstrap 95\% interval & Raw p \\
\midrule\endfirsthead
\multicolumn{5}{l}{\tablename\ \thetable\ (continued)}\\
\toprule
Cohort & Readout & Mean ratio & Bootstrap 95\% interval & Raw p \\
\midrule\endhead
\bottomrule\endfoot
GSE165816 & Fibroblast fraction & 1.425 & 0.873 to 2.298 & 0.438 \\*
GSE165816 & All-cell topic mean & 2.724 & 1.369 to 5.732 & 0.0599 \\
GSE165816 & Within-fibroblast topic mean & 2.889 & 1.059 to 13.966 & 0.042 \\
GSE231643 & Fibroblast fraction & 2.051 & 1.128 to 4.184 & 0.0714 \\
GSE231643 & All-cell topic mean & 2.564 & 1.310 to 8.079 & 0.143 \\*
GSE231643 & Within-fibroblast topic mean & 1.660 & 1.007 to 6.102 & 0.143 \\
\end{longtable}
}

{\small
\begin{longtable}{@{}p{31mm}p{59mm}p{26mm}p{51mm}@{}}
\caption{Historical paired AUC differences for the pooled-fold scoring procedure at each analysis unit. Each contrast subtracts the second leave-one-out readout from the first. Percentile 95\% intervals use 4,000 common-index, arm-stratified bootstrap draws of fixed out-of-fold scores: 9 versus 5 specimens or 7 versus 4 mapped patients. PCA and NMF denote principal component analysis and nonnegative matrix factorization. No sample-unit interval excludes zero; patient-unit Fibroblast topic minus NMF and Module minus NMF intervals do. Because raw PCA/NMF coordinates can change between fitted folds, neither exclusion from zero supports representation superiority. These historical intervals are unadjusted, omit refitting uncertainty and do not describe the revised within-model held-pair estimand.}\label{tab:p1aucdifferences}\\
\toprule
Unit & AUC contrast & Difference & Bootstrap 95\% interval \\
\midrule\endfirsthead
\multicolumn{4}{l}{\tablename\ \thetable\ (continued)}\\
\toprule
Unit & AUC contrast & Difference & Bootstrap 95\% interval \\
\midrule\endhead
\bottomrule\endfoot
14 specimens & Fibroblast topic minus Module & -0.0222 & -0.2889 to 0.2222 \\*
14 specimens & Fibroblast topic minus PCA & 0.2889 & -0.0889 to 0.6667 \\
14 specimens & Fibroblast topic minus NMF & -0.0556 & -0.3556 to 0.2222 \\
14 specimens & Module minus PCA & 0.3111 & 0.0000 to 0.6667 \\
14 specimens & Module minus NMF & -0.0333 & -0.2889 to 0.2000 \\
14 specimens & PCA minus NMF & -0.3444 & -0.7111 to 0.0000 \\
11 patients & Fibroblast topic minus Module & -0.0714 & -0.4286 to 0.2857 \\
11 patients & Fibroblast topic minus PCA & 0.2143 & -0.4643 to 0.9286 \\
11 patients & Fibroblast topic minus NMF & 0.5893 & 0.1429 to 0.9643 \\
11 patients & Module minus PCA & 0.2857 & -0.1786 to 0.7857 \\
11 patients & Module minus NMF & 0.6607 & 0.2857 to 0.9652 \\*
11 patients & PCA minus NMF & 0.3750 & 0.0000 to 0.7857 \\
\end{longtable}
}

{\small
\begin{longtable}{@{}p{42mm}p{35mm}p{48mm}p{21mm}p{19mm}@{}}
\caption{Complete paired anatomical control across 15 topics in 10 author-mapped patients. Differences are mean all-cell loadings in foot minus forearm, after cell-count-weighted collapse of repeated foot specimens. Intervals are percentile 95\% limits from 4,000 paired-patient bootstrap draws. Two-sided p values enumerate all 1,024 sign allocations and retain add-one correction; q values use Benjamini--Hochberg across 15 topics. The TIMP1/COL3A1, COL1A1/COL1A2, CFD/APOD and IGFBP7/TAGLN loadings have q below 0.05. Labels list the two leading decoder genes; they do not assert cell-type identity. Bootstrap intervals and permutation tests need not be dual. This is an anatomical sensitivity control within the discovery source, not a healing endpoint.}\label{tab:p1anatomy}\\
\toprule
Leading genes & Foot minus forearm & Bootstrap 95\% interval & Sign-flip p & BH q \\
\midrule\endfirsthead
\multicolumn{5}{l}{\tablename\ \thetable\ (continued)}\\
\toprule
Leading genes & Foot minus forearm & Bootstrap 95\% interval & Sign-flip p & BH q \\
\midrule\endhead
\bottomrule\endfoot
TIMP1/COL3A1 & 0.0476 & 0.0119 to 0.0916 & 0.0107 & 0.0402 \\*
JUNB/DNAJB1 & 0.0103 & -0.0616 to 0.0793 & 0.799 & 0.956 \\
COL1A1/COL1A2 & 0.0914 & 0.0535 to 0.1376 & 0.00293 & 0.022 \\
IGFBP7/IGFBP5 & 0.0188 & 0.0041 to 0.0361 & 0.0634 & 0.119 \\
TAGLN/ACTA2 & 0.0016 & -0.0200 to 0.0241 & 0.893 & 0.956 \\
KRT14/KRT10 & -0.2491 & -0.4274 to -0.0682 & 0.0361 & 0.0774 \\
CFD/APOD & -0.0362 & -0.0567 to -0.0169 & 0.0107 & 0.0402 \\
DCN/CFD & -0.0159 & -0.0383 to 0.0042 & 0.208 & 0.312 \\
DCD/SAT1 & -0.0001 & -0.0083 to 0.0098 & 0.975 & 0.975 \\
IGFBP7/SPARCL1 & 0.0398 & -0.0015 to 0.0897 & 0.134 & 0.223 \\
FOS/IGFBP7 & -0.0038 & -0.0087 to 0.0025 & 0.245 & 0.334 \\
HLA-DRA/CD74 & 0.0111 & -0.0591 to 0.0984 & 0.854 & 0.956 \\
IGFBP7/TAGLN & 0.0715 & 0.0436 to 0.1039 & 0.00293 & 0.022 \\
FOS/JUNB & 0.0068 & 0.0025 to 0.0118 & 0.0185 & 0.0556 \\*
DNAJB1/FOS & 0.0061 & 0.0020 to 0.0108 & 0.0224 & 0.0561 \\
\end{longtable}
}

{\small
\begin{longtable}{@{}p{33mm}p{23mm}p{33mm}p{33mm}p{43mm}@{}}
\caption{Complete Welch two-one-sided-test sensitivity grids for the all-cell mean fibroblast-associated topic contrast. The symmetric margin is on the loading scale, and both lower and upper one-sided p values must be below 0.05 for the final column to be Yes. Sample rows use 9 versus 5 specimens; patient rows use 7 versus 4 mapped patients with cell-count-weighted repeat collapse. Margins were examined retrospectively and have no externally justified clinical meaning. A Yes entry therefore records a mathematical sensitivity result, not demonstrated clinical equivalence.}\label{tab:p1equivalence}\\
\toprule
Unit & Margin & Lower-test p & Upper-test p & Both below 0.05 \\
\midrule\endfirsthead
\multicolumn{5}{l}{\tablename\ \thetable\ (continued)}\\
\toprule
Unit & Margin & Lower-test p & Upper-test p & Both below 0.05 \\
\midrule\endhead
\bottomrule\endfoot
14 specimens & 0.05 & 0.000198 & 0.413 & No \\*
14 specimens & 0.10 & 5.59e-06 & 0.0077 & Yes \\
14 specimens & 0.20 & 3.66e-08 & 3.31e-06 & Yes \\
14 specimens & 0.30 & 1.14e-09 & 2.6e-08 & Yes \\
14 specimens & 0.50 & 9.95e-12 & 6.75e-11 & Yes \\
11 patients & 0.05 & 0.00172 & 0.159 & No \\
11 patients & 0.10 & 7.56e-05 & 0.00323 & Yes \\
11 patients & 0.20 & 1.09e-06 & 1.02e-05 & Yes \\
11 patients & 0.30 & 6.42e-08 & 3e-07 & Yes \\*
11 patients & 0.50 & 1.45e-09 & 3.72e-09 & Yes \\
\end{longtable}
}

{\small
\begin{longtable}{@{}p{43mm}p{24mm}p{21mm}p{39mm}p{39mm}@{}}
\caption{Patient-level robustness in the discovery cohort. The first row correlates fibroblast mean loading with the high-state fraction after weighting repeated specimen summaries by fibroblast count within 20 patients (25 specimens). Its percentile 95\% interval uses 10,000 patient resamples, seed 217. The second row uses the previously defined all-cell healing contrast in 7 healing and 4 non-healing patients; its original interval remains in the outcome table. Omission ranges give the minimum and maximum after deleting one patient at a time (20 or 11 omissions), not confidence limits. Neither analysis refits the expression model, changes the state cut, adds patients, or provides external validation.}\label{tab:p1robustness}\\
\toprule
Quantity & Unit & Estimate & Patient bootstrap 95\% & Patient-omission range \\
\midrule\endfirsthead
\multicolumn{5}{l}{\tablename\ \thetable\ (continued)}\\
\toprule
Quantity & Unit & Estimate & Patient bootstrap 95\% & Patient-omission range \\
\midrule\endhead
\bottomrule\endfoot
Mean versus high-state fraction & 20 patients & 0.9372 & 0.7760 to 0.9903 & 0.9266 to 0.9643 \\*
Healing minus non-healing & 11 patients & 0.0293 & Not recalculated & 0.0143 to 0.0365 \\
\end{longtable}
}

{\small
\begin{longtable}{@{}p{12mm}p{46mm}p{25mm}p{26mm}p{21mm}p{35mm}@{}}
\caption{Topic-capacity sensitivity after tracking the TIMP1-carrying topic at each K. Gene-based descriptions identify the tracked loading; each fit supplies its own two-Gaussian midpoint cut. Semantic rho uses 25 specimens; the unadjusted healing-arm p value uses 9 versus 5 high-state-fraction summaries. The final column asks whether the maximum fraction occurs in healthy skin. At K=10 the tracked topic is a different IGFBP-rich programme and the counterexample fails. Topics across capacities are not assumed to have identical biological meaning.}\label{tab:p1capacity}\\
\toprule
K & Programme description & Derived cut & Semantic rho & Healing p & Healthy maximum \\
\midrule\endfirsthead
\multicolumn{6}{l}{\tablename\ \thetable\ (continued)}\\
\toprule
K & Programme description & Derived cut & Semantic rho & Healing p & Healthy maximum \\
\midrule\endhead
\bottomrule\endfoot
10 & IGFBP-rich & 0.1454 & 0.9469 & 0.147 & No \\*
15 & CHI3L1/COL3A1-rich & 0.1691 & 0.9316 & 0.109 & Yes \\*
20 & CHI3L1/COL3A1-rich & 0.1695 & 0.9236 & 0.109 & Yes \\
\end{longtable}
}

{\small
\begin{longtable}{@{}p{27mm}p{31mm}p{43mm}p{25mm}p{39mm}@{}}
\caption{All 18 state-cut sensitivity evaluations for the frozen K=15 topic. Spearman rho correlates fibroblast mean loading with the threshold-defined high-state fraction across 25 specimens. The ratio compares mean high-state fractions in 9 healing versus 5 non-healing specimens; p is the unadjusted two-sided Mann--Whitney probability. The last column gives the group containing the largest fraction. The grid includes the fitted cut rounded to 0.1843. Four cuts have p below 0.05; these are dependent sensitivity analyses, not 18 independent confirmations.}\label{tab:p1cuts}\\
\toprule
State cut & Semantic rho & Healing / non-healing & Raw p & Largest fraction \\
\midrule\endfirsthead
\multicolumn{5}{l}{\tablename\ \thetable\ (continued)}\\
\toprule
State cut & Semantic rho & Healing / non-healing & Raw p & Largest fraction \\
\midrule\endhead
\bottomrule\endfoot
0.0800 & 0.9326 & 2.966 & 0.042 & Healthy \\*
0.1000 & 0.9414 & 2.959 & 0.042 & Healthy \\
0.1200 & 0.9091 & 3.011 & 0.042 & Healthy \\
0.1400 & 0.9301 & 2.992 & 0.042 & Healthy \\
0.1600 & 0.9294 & 2.981 & 0.0599 & Healthy \\
0.1800 & 0.9349 & 3.001 & 0.0599 & Healthy \\
0.1843 & 0.9325 & 2.975 & 0.0599 & Healthy \\
0.2000 & 0.9349 & 2.987 & 0.0599 & Healthy \\
0.2200 & 0.9341 & 3.124 & 0.0599 & Healthy \\
0.2400 & 0.9228 & 3.352 & 0.0617 & Healthy \\
0.2600 & 0.9228 & 3.424 & 0.0617 & Healthy \\
0.2800 & 0.9236 & 3.465 & 0.0617 & Healthy \\
0.3000 & 0.9236 & 3.499 & 0.0617 & Healthy \\
0.3200 & 0.9479 & 3.741 & 0.0617 & Healthy \\
0.3400 & 0.9463 & 3.699 & 0.0617 & Healthy \\
0.3600 & 0.9401 & 4.040 & 0.0608 & Healthy \\
0.3800 & 0.9182 & 3.945 & 0.0608 & Healthy \\*
0.4000 & 0.9182 & 3.889 & 0.0608 & Healthy \\
\end{longtable}
}

{\small
\begin{longtable}{@{}p{63mm}p{55mm}p{49mm}@{}}
\caption{Exact mean decomposition on the saved 19,410 discovery fibroblasts in 25 specimens. The cut is 0.184295848; empirical bin means are distinct from fitted Gaussian-component means. The constant-state approximation uses the cell-pooled low/high means, and its residual is the observed specimen mean minus that approximation. RMS denotes root-mean-square residual with equal specimen weights. Ranges exclude undefined means for empty bins; zero-weight terms contribute zero to the identity. Both exact identities have maximum numerical discrepancy below $10^{-15}$. Values are on the zero-to-one loading scale, apart from the count row. They describe the frozen sample, without confidence intervals, a test of invariant state intensity or causal variance attribution.}\label{tab:p1math}\\
\toprule
Quantity & Value or range & Denominator \\
\midrule\endfirsthead
\multicolumn{3}{l}{\tablename\ \thetable\ (continued)}\\
\toprule
Quantity & Value or range & Denominator \\
\midrule\endhead
\bottomrule\endfoot
Cell-pooled low-bin mean & 0.02380 & All 19,410 fibroblasts \\*
Cell-pooled high-bin mean & 0.51301 & All 19,410 fibroblasts \\
Specimen low-bin means & 0.01204 to 0.10213 & 25 nonempty low bins \\
Specimen high-bin means & 0.22645 to 0.57777 & 16 nonempty high bins \\
Specimens with an empty bin & 9 & 9 of 25; high bin empty \\
Within-state residual range & -0.02434 to 0.04377 & All 25 specimens \\*
Within-state residual RMS & 0.01575 & Equal specimen weights \\
\end{longtable}
}

{\small
\begin{longtable}{@{}p{100mm}p{67mm}@{}}
\caption{Independent arithmetic check of the all-cell mean-loading contrast in 11 patients (7 healing, 4 non-healing). Exhaustive allocation includes the observed labels and numerical ties: 90/330 allocations are at least as extreme in absolute mean difference, whereas the retained add-one convention gives 91/331. These probabilities assume observational label exchangeability and do not identify a causal effect. Welch quantities use sample variances within the two outcome groups. Standard error, interval and margin are on the loading scale; degrees of freedom and probabilities are dimensionless. A margin strictly greater than the infimum would be needed to reject both one-sided tests at 0.05. The data-derived boundary is not a clinical equivalence margin.}\label{tab:p1exact}\\
\toprule
Quantity & Value \\
\midrule\endfirsthead
\multicolumn{2}{l}{\tablename\ \thetable\ (continued)}\\
\toprule
Quantity & Value \\
\midrule\endhead
\bottomrule\endfoot
All label allocations & 330 \\*
At least as extreme as observed & 90 \\
Exhaustive two-sided probability & 0.272727 \\
Recorded add-one probability & 0.274924 \\
Welch standard error & 0.019446 \\
Welch degrees of freedom & 8.113833 \\
Two-sided 90\% Welch interval & -0.006816 to 0.065375 \\*
Equivalence-margin infimum & 0.065375 \\
\end{longtable}
}

{\small
\begin{longtable}{@{}p{15mm}p{25mm}p{18mm}p{23mm}p{27mm}p{25mm}p{27mm}@{}}
\caption{Observed-cell extension: mapped patient-level tissue measurements from 25 selected foot specimens. Patient labels are stable pseudonyms within this extension. The 20 rows include 9 non-diabetic participants, who do not enter healing tests, and 11 DFU patients (7 healed and 4 not healed). Repeated specimens are pooled by fibroblast count for the mean and high-state fraction, and by all-cell count for fibroblast composition. High state uses the unchanged discovery midpoint. Specimens and fibroblasts are counts; the three readouts are unitless fractions or loadings on zero-to-one scales.}\label{tab:p1biologicalpatients}\\
\toprule
Patient & Group & Specimens & Fibroblasts & Fibroblast mean & High fraction & Fibroblast fraction \\
\midrule\endfirsthead
\multicolumn{7}{l}{\tablename\ \thetable\ (continued)}\\
\toprule
Patient & Group & Specimens & Fibroblasts & Fibroblast mean & High fraction & Fibroblast fraction \\
\midrule\endhead
\bottomrule\endfoot
01 & Not healed & 1 & 221 & 0.2496 & 0.5113 & 0.0970 \\*
02 & Healed & 2 & 2751 & 0.3473 & 0.5863 & 0.3670 \\
03 & Healed & 1 & 1238 & 0.0274 & 0.0226 & 0.3133 \\
04 & Healed & 1 & 356 & 0.3521 & 0.6994 & 0.1582 \\
05 & Not healed & 1 & 843 & 0.0125 & 0.0000 & 0.1965 \\
06 & Not healed & 2 & 1552 & 0.0179 & 0.0032 & 0.2021 \\
07 & Not healed & 1 & 872 & 0.0547 & 0.0837 & 0.3515 \\
08 & Healed & 2 & 1605 & 0.1648 & 0.2960 & 0.3829 \\
09 & Healed & 1 & 363 & 0.2087 & 0.3719 & 0.1240 \\
10 & Healed & 1 & 170 & 0.0705 & 0.0824 & 0.0925 \\
11 & Healed & 1 & 1821 & 0.0393 & 0.0483 & 0.5397 \\
12 & Non-diabetic & 1 & 1020 & 0.4663 & 0.8412 & 0.3491 \\
13 & Non-diabetic & 1 & 215 & 0.0295 & 0.0186 & 0.0866 \\
14 & Non-diabetic & 2 & 1921 & 0.0137 & 0.0000 & 0.2509 \\
15 & Non-diabetic & 1 & 1152 & 0.0120 & 0.0000 & 0.3904 \\
16 & Non-diabetic & 2 & 1562 & 0.0153 & 0.0000 & 0.1873 \\
17 & Non-diabetic & 1 & 153 & 0.0136 & 0.0000 & 0.1939 \\
18 & Non-diabetic & 1 & 974 & 0.0131 & 0.0010 & 0.1804 \\
19 & Non-diabetic & 1 & 425 & 0.0135 & 0.0000 & 0.2941 \\*
20 & Non-diabetic & 1 & 196 & 0.0205 & 0.0000 & 0.1142 \\
\end{longtable}
}

{\small
\begin{longtable}{@{}p{34mm}p{34mm}p{23mm}p{36mm}p{17mm}p{17mm}@{}}
\caption{Patient-unit extension of the three fibroblast readouts. Differences are healed minus not healed in the same 7 and 4 DFU patients. Percentile intervals use 10,000 within-arm patient resamples, seed 17. Exact two-sided probabilities enumerate all 330 label allocations including ties, without add-one correction; q values apply Benjamini--Hochberg adjustment to the three readouts within each weighting variant. These are retrospective observational comparisons without a clinical equivalence margin. Every DFU specimen already has at least 100 fibroblasts, so the 100-cell eligibility variant reproduces the pooled-cell comparison exactly and duplicate rows are omitted. Repeated-cell and specimen weighting affect only aggregation within a patient.}\label{tab:p1biologicalcontrasts}\\
\toprule
Aggregation & Readout & Difference & Bootstrap 95\% interval & Exact p & BH q \\
\midrule\endfirsthead
\multicolumn{6}{l}{\tablename\ \thetable\ (continued)}\\
\toprule
Aggregation & Readout & Difference & Bootstrap 95\% interval & Exact p & BH q \\
\midrule\endhead
\bottomrule\endfoot
Pooled cells & Fibroblast mean & 0.0892 & -0.0486 to 0.2177 & 0.3121 & 0.4697 \\*
Pooled cells & High-state fraction & 0.1514 & -0.1331 to 0.4153 & 0.3697 & 0.4697 \\
Pooled cells & Fibroblast composition & 0.0708 & -0.0708 to 0.2146 & 0.4697 & 0.4697 \\
Equal specimens & Fibroblast mean & 0.0974 & -0.0421 to 0.2260 & 0.2788 & 0.4848 \\
Equal specimens & High-state fraction & 0.1683 & -0.1170 to 0.4310 & 0.3424 & 0.4848 \\*
Equal specimens & Fibroblast composition & 0.0679 & -0.0724 to 0.2114 & 0.4848 & 0.4848 \\
\end{longtable}
}

{\small
\begin{longtable}{@{}p{28mm}p{24mm}p{32mm}p{33mm}p{47mm}@{}}
\caption{Patient-level threshold sensitivity in the same 7 healed and 4 non-healed DFU patients. Each patient has equal group weight; recovered fibroblasts are pooled within patient. Differences are healed minus non-healed high-state fractions. Exact two-sided p values enumerate all 330 patient-label allocations, including ties, without add-one correction. The support-wide p compares each absolute contrast with the maximum absolute patient-group CDF difference over every distinct observed loading in each allocation; it accounts for searching across that complete support under the exchangeability null. The threshold-free maximum is 0.3045 with exact p=0.1848. Original denotes the unchanged fitted midpoint. These retrospective diagnostics are not independent validation.}\label{tab:p1thresholdpatient}\\
\toprule
Threshold & Original & Difference & Exact p & Support-wide p \\
\midrule\endfirsthead
\multicolumn{5}{l}{\tablename\ \thetable\ (continued)}\\
\toprule
Threshold & Original & Difference & Exact p & Support-wide p \\
\midrule\endhead
\bottomrule\endfoot
0.0500 & False & 0.2485 & 0.2242 & 0.3182 \\*
0.1000 & False & 0.1958 & 0.3030 & 0.4515 \\
0.1500 & False & 0.1682 & 0.3424 & 0.5606 \\
0.1843 & True & 0.1514 & 0.3697 & 0.6152 \\
0.2000 & False & 0.1445 & 0.3727 & 0.6424 \\
0.3000 & False & 0.1300 & 0.3727 & 0.7182 \\
0.4000 & False & 0.1029 & 0.3909 & 0.8061 \\
0.5000 & False & 0.0819 & 0.3667 & 0.8758 \\*
0.6000 & False & 0.0570 & 0.3091 & 0.9485 \\
\end{longtable}
}

{\small
\begin{longtable}{@{}p{30mm}p{31mm}p{23mm}p{32mm}p{32mm}p{16mm}@{}}
\caption{Cell-count and repeated-specimen sensitivity at the fixed original threshold. Equal-cell rows use 250 draws of 100 fibroblasts without replacement from each DFU specimen, seed 29, followed by equal specimen weighting within each patient. The other rows enumerate every one-specimen-per-patient choice across the three repeated patients. Every comparison retains 7 healed and 4 non-healed patients and recomputes the exact 330-allocation probability. Difference ranges are healed minus non-healed; positive counts cover the listed number of comparisons. Neither range is a confidence interval. Composition is omitted from equal-fibroblast sampling because that selected subset no longer estimates the original all-cell denominator.}\label{tab:p1samplingsensitivity}\\
\toprule
Sensitivity & Readout & Comparisons & Difference range & Exact p range & Positive \\
\midrule\endfirsthead
\multicolumn{6}{l}{\tablename\ \thetable\ (continued)}\\
\toprule
Sensitivity & Readout & Comparisons & Difference range & Exact p range & Positive \\
\midrule\endhead
\bottomrule\endfoot
100 cells per specimen & Fibroblast mean & 250 & 0.0749 to 0.1233 & 0.1273 to 0.3879 & 250/250 \\*
100 cells per specimen & High-state fraction & 250 & 0.1098 to 0.2229 & 0.1515 to 0.5667 & 250/250 \\
One specimen per patient & Fibroblast mean & 8 & 0.0777 to 0.1170 & 0.2333 to 0.3303 & 8/8 \\
One specimen per patient & High-state fraction & 8 & 0.1284 to 0.2081 & 0.2667 to 0.4394 & 8/8 \\*
One specimen per patient & Fibroblast composition & 8 & 0.0440 to 0.0918 & 0.3909 to 0.6515 & 8/8 \\
\end{longtable}
}

{\small
\begin{longtable}{@{}p{73mm}p{47mm}p{47mm}@{}}
\caption{Raw-count column-shuffle control on 6,000 cells, with 240 cells sampled from each of 25 specimens. Every one of 23,572 gene-value multisets is preserved exactly before normalization. Both branches use the same frozen encoder and deterministic projection. HHI is the sum of squared topic proportions; lower HHI and higher perplexity describe reduced concentration. Library sizes and detected-gene counts use all source genes. Detection fraction divides the detected-gene count by 23,572. Means are preserved by column marginals, but row-wise variation narrows. These descriptive values neither isolate biological co-expression from technical row structure nor show rejection by the historical integrity gate.}\label{tab:p1rawcontrol}\\
\toprule
Quantity & Original counts & Shuffled counts \\
\midrule\endfirsthead
\multicolumn{3}{l}{\tablename\ \thetable\ (continued)}\\
\toprule
Quantity & Original counts & Shuffled counts \\
\midrule\endhead
\bottomrule\endfoot
Mean topic perplexity & 3.4383 & 11.2306 \\*
Mean HHI & 0.5740 & 0.1467 \\
Library size, mean & 9275.2065 & 9275.2065 \\
Library size, SD & 10907.9970 & 1500.9520 \\
Detected genes, mean & 2312.1958 & 2312.1958 \\
Detected genes, SD & 1367.1079 & 38.9770 \\
Detection fraction, mean & 0.0981 & 0.0981 \\*
Detection fraction, SD & 0.0580 & 0.0017 \\
\end{longtable}
}

{\small
\begin{longtable}{@{}p{43mm}p{51mm}p{25mm}p{24mm}p{24mm}@{}}
\caption{Marker-reference sensitivity in 11 DFU patients. Re-estimating marker means and SDs within each specimen changes 7,044 of 76,461 cell labels (9.21\%); fibroblast calls change from 19,410 to 19,711, with 1,119 lost and 1,420 gained. The maximum specimen changes are 0.0371 for the fibroblast mean and 0.0695 for high-state fraction. Listed differences are healed minus non-healed after count-weighted patient aggregation; exact probabilities enumerate 330 allocations, and BH q adjusts the three readouts within each gate. Applying the frozen pooled reference per specimen reproduces all original labels and is omitted as a duplicate implementation check. No independent cell-identity truth set selects the correct gate.}\label{tab:p1gatecontrol}\\
\toprule
Marker reference & Readout & Difference & Exact p & BH q \\
\midrule\endfirsthead
\multicolumn{5}{l}{\tablename\ \thetable\ (continued)}\\
\toprule
Marker reference & Readout & Difference & Exact p & BH q \\
\midrule\endhead
\bottomrule\endfoot
Pooled reference & Fibroblast mean & 0.0892 & 0.3121 & 0.4697 \\*
Pooled reference & High-state fraction & 0.1514 & 0.3697 & 0.4697 \\
Pooled reference & Fibroblast composition & 0.0708 & 0.4697 & 0.4697 \\
Per-specimen reference & Fibroblast mean & 0.0887 & 0.3182 & 0.4970 \\
Per-specimen reference & High-state fraction & 0.1558 & 0.3667 & 0.4970 \\*
Per-specimen reference & Fibroblast composition & 0.0553 & 0.4970 & 0.4970 \\
\end{longtable}
}

{\small
\begin{longtable}{@{}p{54mm}p{28mm}p{54mm}p{31mm}@{}}
\caption{Secondary DFU-only specimen-label calibration. All 11 healthy specimens are excluded before enumerating the 2,002 allocations of 9 healed and 5 non-healed specimen labels. Differences are equal-specimen arm means, healed minus non-healed. The central 95\% null interval describes the allocation distribution and is not an effect confidence interval. Exact two-sided probabilities include the observed allocation and ties without add-one correction. Repeated specimens share patients, so the 11-patient comparisons remain primary.}\label{tab:p1dfunull}\\
\toprule
Readout & Difference & Central null 95\% & Exact p \\
\midrule\endfirsthead
\multicolumn{4}{l}{\tablename\ \thetable\ (continued)}\\
\toprule
Readout & Difference & Central null 95\% & Exact p \\
\midrule\endhead
\bottomrule\endfoot
Fibroblast mean & 0.1335 & -0.1652 to 0.1623 & 0.0894 \\*
High-state fraction & 0.2380 & -0.3167 to 0.3222 & 0.1199 \\*
Fibroblast composition & 0.0888 & -0.1540 to 0.1460 & 0.2697 \\
\end{longtable}
}

{\small
\begin{longtable}{@{}p{24mm}p{25mm}p{25mm}p{35mm}p{29mm}p{23mm}@{}}
\caption{Disagreement between five repeated-specimen pairs from five patients, including three DFU patients. G labels retain source specimen identifiers. Mean and high-fraction differences are absolute, using fibroblast-restricted summaries and the unchanged threshold. The split-half column lists the 97.5th percentile of absolute mean differences from 250 random disjoint fibroblast partitions within each specimen, in pair order. The empirical CDF supremum compares the two full loading distributions. Cell partitions describe conditional finite-cell variability; they are not biological replicates, confidence intervals, a patient-level null, an intraclass correlation or clinical reliability estimates.}\label{tab:p1repeats}\\
\toprule
Specimens & Group & Mean difference & Split-half limits & High-fraction diff. & CDF supremum \\
\midrule\endfirsthead
\multicolumn{6}{l}{\tablename\ \thetable\ (continued)}\\
\toprule
Specimens & Group & Mean difference & Split-half limits & High-fraction diff. & CDF supremum \\
\midrule\endhead
\bottomrule\endfoot
G23 / G4 & Healed & 0.0096 & 0.0309 / 0.0445 & 0.0025 & 0.0576 \\*
G33 / G34 & Not healed & 0.0042 & 0.0017 / 0.0041 & 0.0053 & 0.1251 \\
G7 / G8 & Healed & 0.2583 & 0.0553 / 0.0185 & 0.5458 & 0.5789 \\
G31 / G32 & Non-diabetic & 0.0002 & 0.0007 / 0.0009 & 0.0000 & 0.0337 \\*
G10 / G24 & Non-diabetic & 0.0019 & 0.0010 / 0.0013 & 0.0000 & 0.1865 \\
\end{longtable}
}

{\small
\begin{longtable}{@{}p{37mm}p{18mm}p{18mm}p{20mm}p{33mm}p{36mm}@{}}
\caption{Revised same-model held-pair discrimination in 11 mapped DFU patients (7 healing, 4 non-healing). Each of the 28 folds holds one patient from each arm out, fits on the other nine and scores both with one unchanged rule. Wins favour the healing patient; an exact tie earns one-half credit. Discrimination is total credit divided by 28, not a pooled-fold raw-coordinate AUC. Topic and module orientation and PCA/NMF fitting and component selection use training patients. NMF basis normalization and training-SD-standardized component selection remove positive scaling ambiguity. The uniform NMF cap was raised from 1,000 to 10,000 after training-convergence warnings; no final fit warned. All pairs share 11 patients, and discovery expression informed the fixed panel and encoder. No confidence intervals, p values, independent validation or representation ranking are claimed.}\label{tab:p1pairauc}\\
\toprule
Readout & Wins & Ties & Losses & Pair credit & Discrimination \\
\midrule\endfirsthead
\multicolumn{6}{l}{\tablename\ \thetable\ (continued)}\\
\toprule
Readout & Wins & Ties & Losses & Pair credit & Discrimination \\
\midrule\endhead
\bottomrule\endfoot
Frozen topic & 20 & 0 & 8 & 20.0 / 28 & 0.7143 \\*
Module & 22 & 0 & 6 & 22.0 / 28 & 0.7857 \\
PCA & 14 & 0 & 14 & 14.0 / 28 & 0.5000 \\*
NMF & 5 & 3 & 20 & 6.5 / 28 & 0.2321 \\
\end{longtable}
}

{\small
\begin{longtable}{@{}p{25mm}p{27mm}p{25mm}p{31mm}p{30mm}p{27mm}@{}}
\caption{Full released-count depth audit of all 81,602 pre-analysis-QC columns from 25 foot-skin specimens. Baseline reproduces 76,461 retained cells and 19,410 fibroblast labels. Each nonzero count is independently binomial-thinned at the stated probability; each rate starts from original counts and uses all three fixed computational seeds. QC is at least 200 detected genes and mitochondrial fraction at most 0.20. Marker means/SDs, forced-winner gate, encoder and high-state cut remain frozen. Lost and gained cells are relative to baseline QC, not a net count. Lineage switches divide changed labels by cells surviving both QC passes. Zero-marker cells pass QC but detect none of the annotation markers; their forced labels are not independently validated. Counts are cells, not independent subjects; seeds provide no biological replication.}\label{tab:p1fulldepth}\\
\toprule
Count retention & Seed & QC cells & Lost / gained & Lineage switches & Zero-marker cells \\
\midrule\endfirsthead
\multicolumn{6}{l}{\tablename\ \thetable\ (continued)}\\
\toprule
Count retention & Seed & QC cells & Lost / gained & Lineage switches & Zero-marker cells \\
\midrule\endhead
\bottomrule\endfoot
100\% & Baseline & 76461 & 0 / 0 & 0.00\% & 92 \\*
75\% & 20260925 & 76375 & 127 / 41 & 3.18\% & 165 \\
75\% & 20260926 & 76387 & 113 / 39 & 3.30\% & 169 \\
75\% & 20260927 & 76393 & 117 / 49 & 3.29\% & 168 \\
50\% & 20260925 & 75738 & 793 / 70 & 6.23\% & 384 \\
50\% & 20260926 & 75726 & 799 / 64 & 6.25\% & 363 \\
50\% & 20260927 & 75724 & 810 / 73 & 6.34\% & 386 \\
25\% & 20260925 & 72145 & 4419 / 103 & 10.29\% & 770 \\
25\% & 20260926 & 72197 & 4367 / 103 & 10.48\% & 842 \\
25\% & 20260927 & 72145 & 4404 / 88 & 10.45\% & 815 \\
10\% & 20260925 & 61397 & 15175 / 111 & 15.36\% & 1429 \\
10\% & 20260926 & 61366 & 15210 / 115 & 15.40\% & 1462 \\*
10\% & 20260927 & 61395 & 15178 / 112 & 15.31\% & 1410 \\
\end{longtable}
}

{\small
\begin{longtable}{@{}p{25mm}p{48mm}p{26mm}p{35mm}p{33mm}@{}}
\caption{Patient contrasts after rerunning QC and the frozen-reference gate on each depth perturbation. Differences are healed minus non-healed, with seven and four author-mapped DFU patients. Repeated specimens are pooled by the appropriate all-cell or fibroblast count. Fibroblast composition uses all current QC survivors; high state uses the unchanged discovery cut. The exact two-sided probability enumerates all 330 patient-label allocations, including ties, without add-one correction. Baseline has one run; other rows give the minimum and maximum over three fixed thinning seeds. These are computational ranges, not confidence intervals, multiplicity-adjusted discoveries, equivalence bounds or biological replication. Zero-panel cells remain explicitly missing in conditional branches rather than receiving scores from an all-zero vector.}\label{tab:p1depthoutcomes}\\
\toprule
Count retention & Readout & Estimable runs & Difference range & Exact p range \\
\midrule\endfirsthead
\multicolumn{5}{l}{\tablename\ \thetable\ (continued)}\\
\toprule
Count retention & Readout & Estimable runs & Difference range & Exact p range \\
\midrule\endhead
\bottomrule\endfoot
100\% & All-cell mean & 1 / 1 & 0.0293 to 0.0293 & 0.2727 to 0.2727 \\*
100\% & Fibroblast mean & 1 / 1 & 0.0892 to 0.0892 & 0.3121 to 0.3121 \\
100\% & High-state fraction & 1 / 1 & 0.1514 to 0.1514 & 0.3697 to 0.3697 \\
100\% & Fibroblast composition & 1 / 1 & 0.0708 to 0.0708 & 0.4697 to 0.4697 \\
75\% & All-cell mean & 3 / 3 & 0.0291 to 0.0292 & 0.2758 to 0.2818 \\
75\% & Fibroblast mean & 3 / 3 & 0.0876 to 0.0899 & 0.2970 to 0.3182 \\
75\% & High-state fraction & 3 / 3 & 0.1451 to 0.1518 & 0.3545 to 0.3788 \\
75\% & Fibroblast composition & 3 / 3 & 0.0687 to 0.0706 & 0.4727 to 0.4788 \\
50\% & All-cell mean & 3 / 3 & 0.0289 to 0.0291 & 0.2818 to 0.2848 \\
50\% & Fibroblast mean & 3 / 3 & 0.0880 to 0.0913 & 0.2848 to 0.3152 \\
50\% & High-state fraction & 3 / 3 & 0.1485 to 0.1620 & 0.3364 to 0.3727 \\
50\% & Fibroblast composition & 3 / 3 & 0.0683 to 0.0698 & 0.4758 to 0.4848 \\
25\% & All-cell mean & 3 / 3 & 0.0292 to 0.0293 & 0.2788 to 0.2879 \\
25\% & Fibroblast mean & 3 / 3 & 0.0883 to 0.0902 & 0.2758 to 0.2758 \\
25\% & High-state fraction & 3 / 3 & 0.1526 to 0.1634 & 0.3303 to 0.3424 \\
25\% & Fibroblast composition & 3 / 3 & 0.0737 to 0.0759 & 0.4455 to 0.4727 \\
10\% & All-cell mean & 3 / 3 & 0.0288 to 0.0296 & 0.3303 to 0.3455 \\
10\% & Fibroblast mean & 3 / 3 & 0.0879 to 0.0918 & 0.1970 to 0.2212 \\
10\% & High-state fraction & 3 / 3 & 0.1645 to 0.1687 & 0.2515 to 0.3000 \\*
10\% & Fibroblast composition & 3 / 3 & 0.0845 to 0.0897 & 0.3636 to 0.4091 \\
\end{longtable}
}
}

\end{document}
